\documentclass[11pt]{article}

\usepackage[a4paper,margin=24mm]{geometry}
\usepackage[T1]{fontenc}
\usepackage{lmodern}
\usepackage{microtype}
\usepackage{amsmath,amssymb,mathtools,amsthm}
\usepackage{booktabs,tabularx,array,longtable}
\usepackage{graphicx}
\usepackage{xcolor}
\usepackage{enumitem}
\usepackage{xurl}
\usepackage[hidelinks]{hyperref}
\usepackage[backend=biber,style=numeric,sorting=none,maxbibnames=10]{biblatex}
\addbibresource{references.bib}

\setlength{\emergencystretch}{3em}
\renewcommand{\tabularxcolumn}[1]{>{\raggedright\arraybackslash}p{#1}}
\definecolor{statusblue}{RGB}{31,78,121}
\definecolor{statusgray}{RGB}{244,246,248}
\newtheorem{theorem}{Theorem}
\newtheorem{proposition}[theorem]{Proposition}
\newtheorem{corollary}[theorem]{Corollary}
\newcommand{\norm}[1]{\left\lVert#1\right\rVert}
\newcommand{\C}{\mathbb C}
\newcommand{\R}{\mathbb R}
\newcommand{\Hsrc}{H_{\mathrm W}}
\newcommand{\Hred}{H_{\mathrm r}}
\newcommand{\kvec}{\mathbf k}
\newcommand{\qvec}{\mathbf q}
\newcommand{\Rvec}{\mathbf R}
\newcommand{\statusbox}[1]{%
  \begin{center}
  \fcolorbox{statusblue}{statusgray}{\parbox{0.92\textwidth}{#1}}
  \end{center}}
\newcommand{\resultfigure}[4]{%
  \begin{figure}[tbp]
    \centering
    \includegraphics[width=0.96\textwidth,height=0.68\textheight,keepaspectratio]{#1}
    \caption{#2 #3}
    \label{#4}
  \end{figure}}

\title{\textbf{Band Equivalence, Hopping Locality, and Internal Fidelity\\
of Reduced Wannier Hamiltonians}\\[1mm]
\large From Exact Rank-Two Geometry to a DFT-Derived Rank-10 Silicon Case Study}
\author{Eugene Joseph M. Ragasa\\
Department of Physics, De La Salle University, Manila, Philippines\\
\texttt{eugene.ragasa@dlsu.edu.ph}}
\date{}

\begin{document}
\maketitle

\begin{abstract}
Band energies are invariant under momentum-dependent unitary similarity, whereas a
Hamiltonian's matrix representation and real-space hopping locality are not.  We combine
an exact rank-two theory of this distinction with a source-relative case study of a
DFT-derived rank-10 scalar-relativistic silicon Wannier Hamiltonian.  In the controlled
rank-two family, spectral and globally aligned operator losses have an exact
common-curvature geometry.  For
$U_a(q)=R(a\sin2\pi q)$, the minimum restricted-alignment threshold is
$\sqrt{\kappa[1-J_0(2a)^2]}$; for a flat traceless parent, its square equals the
normalized Fourier power transferred from the onsite block to nonzero displacement.
A dispersive development calculation preserves complete bands while changing the first
sampled cutoff meeting a $10^{-8}E_{\rm ref}$ band-error target from 4 to 10.

The silicon study uses source files whose provenance, byte identities, native optimizer
completion, Hermiticity, pair closure, and reconstruction were verified; this
``authenticated'' status is not material validation.  A frozen internal gauge attack
preserves the complete represented spectrum to $3.20\times10^{-14}$~eV while changing
restricted alignment and finite-block locality.  Frozen shell-reduced models then show
that the 408-point fitting inventory understates exact Parseval operator loss and
dense-mesh spectral loss, and that global spectral RMS does not control conduction
curvature or the nearly degenerate valence manifold.  A polar-transported three-band
quadratic Hamiltonian is therefore used instead of independently ordered hole tensors.

The source fails separate DFT-parent gap and electron-mass targets, and the silicon
results remain producer-checked development evidence pending independent reconstruction.
Within those boundaries, the results establish that spectrum, represented operator,
hopping locality, interpolation fidelity, and band-edge fidelity are distinct reduction
objectives; near-full native support is required to preserve them simultaneously for the
studied source.
\end{abstract}

\noindent\textbf{Keywords:} band equivalence; effective mass; hopping locality;
model reduction; silicon; Wannier Hamiltonian

\section{Introduction}

First-principles electronic-structure calculations produce operators and states in
representations designed for numerical convergence rather than for compact real-space
modeling.  Plane waves are systematic and broadly transferable, but they are expensive
for repeated interpolation, large supercells, local perturbations, and reduced-device
calculations.  Wannier representations replace a selected Bloch subspace by localized
coordinates and expose a Fourier hierarchy of matrix-valued hopping blocks
\cite{marzari1997,souza2001,marzari2012,pizzi2020,yates2007}.  A further reduced model
may then retain fewer blocks, impose symmetry patterns, or fit a smaller parameter
family.

The apparent objective---``reproduce the bands with a shorter Hamiltonian''---contains
several different mathematical tasks.  Band energies constrain eigenvalues.  An
operator comparison constrains matrices after identifying their state spaces.  A
locality comparison constrains Fourier coefficients in a declared basis and gauge.
Interpolation asks how a frozen model behaves away from the points used to characterize
it.  Effective masses and valley locations depend on local derivatives and extrema.
These requirements are related, but none is interchangeable with all the others.

The gauge freedom of a retained Bloch manifold makes the distinction unavoidable.  If
$H(\kvec)$ is transformed by a momentum-dependent unitary $U(\kvec)$,
\begin{equation}
 \widetilde H(\kvec)=U(\kvec)^\dagger H(\kvec)U(\kvec),
 \label{eq:intro-gauge}
\end{equation}
then the complete eigenvalue set is unchanged at every momentum.  The matrix field and
its Fourier coefficients generally change.  A momentum-independent unitary merely
conjugates each real-space block; a momentum-dependent unitary convolves displacement
indices and can redistribute weight among displacement shells, including toward longer
range for particular transformations and locality measures.  No universal monotonic
increase is implied.  Band equivalence therefore does not determine hopping locality.

This paper develops that statement through two complementary but logically distinct
levels.  The first is a controlled theoretical result and synthetic benchmark.  The
second is a source-specific computational case study; it is not an empirical proof of a
universal rank-$N$ extension of the rank-two identities.

\begin{enumerate}[leftmargin=2em]
 \item A controlled one-dimensional rank-two problem supplies exact loss geometry,
 Bessel-function threshold laws, a residual--Fourier-power identity, and a dispersive
 fixed-accuracy range-inflation calculation.
 \item A DFT-derived rank-10 silicon Wannier Hamiltonian tests the same distinctions in
 a higher-rank, three-dimensional represented field with realistic matrix structure,
 native Wigner--Seitz blocks, dense interpolation, conduction valleys, and a degenerate
 valence edge.
\end{enumerate}

The silicon calculation is not the ``validation'' of the synthetic theorem in the sense
of proving production material accuracy.  It is the next rung of a validation ladder:
the exact mechanism is first isolated in a known model, then deployed on an authenticated
material-derived operator while keeping parent-model, representation, conversion,
truncation, fitting, and feature errors separate.

The main questions are:

\begin{enumerate}[leftmargin=2em]
 \item How can complete spectra remain identical while restricted alignment and
 real-space locality change?
 \item What exact relationship connects irreducible matrix mismatch to induced hopping
 weight in a controlled rank-two family?
 \item How does representation-dependent locality affect the finite range needed at
 equal accuracy?
 \item How much native support is required to compress the rank-10 silicon Wannier
 source by literal shell deletion or bounded shell refitting?
 \item Do errors measured on a finite 408-point inventory predict exact integrated
 operator error, dense interpolation error, and band-edge fidelity?
 \item What is the appropriate hole-sector object when the non-SOC valence edge is a
 nearly degenerate three-band manifold?
\end{enumerate}

The answers are deliberately bounded.  This article is a long-form synthesis of the
exact rank-two framework and the completed rank-10 silicon calculation program.
Analytic statements are identified separately from sampled numerical observations.  The
active rank-two numerical headline remains development or partial unsealed evidence
according to its calculation records, and the non-flat numerical results are sampled on
declared inventories.  The silicon gauge, shell, fidelity, and band-edge studies are
authenticated producer-checked development analyses that have not been independently
reconstructed.  They are internal to a source that fails separate DFT-parent adequacy
targets and do not validate the source as an accurate silicon Hamiltonian.  These
boundaries are not qualifications added at the end; they define the scientific question
from the beginning.  Preparation of this working manuscript included one separately
authorized fixed-density Quantum ESPRESSO parent band-path calculation after the
reduction models had been frozen; it changed no source construction, fit, selection, or
threshold.  No Wannier90 execution was performed, and manuscript preparation does not
itself authorize submission, publication, or release.

\section{Background and conceptual framework}

\subsection{Physical model, represented operator, and finite matrix}

A physical electronic-structure model, an operator on a retained state space, and a
finite matrix representation are different objects.  The Kohn--Sham calculation defines
an approximation to the physical parent under a chosen functional, pseudopotential,
geometry, and numerical setting \cite{hohenberg1964,kohn1965,perdew1996,hamann2013}.
Wannierization selects and represents a finite subspace.  A reduced hopping model then
approximates that represented operator.  Errors introduced at these stages must not be
collapsed into one number unless a reviewed contract explicitly defines the combination.

For a finite-rank periodic representation,
\begin{equation}
 H(\kvec)=\sum_{\Rvec} e^{i\kvec\cdot\Rvec}T(\Rvec),
 \label{eq:fourier-main}
\end{equation}
where $T(\Rvec)$ is a matrix block associated with a lattice translation.  Equation
\eqref{eq:fourier-main} is meaningful only after the lattice basis, reciprocal
coordinates, representative convention, degeneracies, units, basis ordering, and
energy reference are fixed.

The silicon program uses two Fourier routes.  Native Wannier90 HR data use
Wigner--Seitz representatives and degeneracy factors.  The generic gauge study samples a
matrix field on a centered finite grid, performs a finite-cyclic transform, and then
converts explicitly to ordinary lattice representatives.  These routes answer related
questions but are not interchangeable.  Their conversion and interpolation consequences
are retained as separate channels.

\subsection{Spectral reduction and finite representation extent}

The distinction between representation and reduction can be seen most cleanly in two
retained one-dimensional calibration problems.  In the particle-in-a-box calculation,
a continuum differential operator, its finite-difference matrix, the computed matrix
eigenpairs, and a low-energy spectral compression are separate objects.  An eigensolver
may solve the finite matrix to high accuracy while that matrix still approximates the
continuum spectrum poorly.  If the retained vectors are exact eigenvectors of the finite
matrix, however, compression to their span is algebraically exact for that represented
matrix.  With
\begin{equation}
 P_M=\sum_{n=1}^{M}\lvert u_n\rangle\langle u_n\rvert,
 \qquad Q_M=I-P_M,
 \qquad H_M=P_MHP_M,
 \label{eq:spectral-reduction-projectors}
\end{equation}
one has $Q_MHP_M=0$ for an invariant spectral subspace
\cite{shankar1994,teschl2014,hornJohnson2013}.  This identity establishes closure of the
reduction; it does not establish that the finite parent represents the continuum or a
material accurately.  For a non-invariant trial space, leakage $Q_MHP_M$ becomes a
separate error channel.

The harmonic-oscillator calibration introduces a second limitation: finite spatial
extent.  A real-line operator represented on a finite Dirichlet interval has changed its
domain and boundary realization before the grid is refined.  Decreasing the grid spacing
at fixed interval width resolves that finite-box problem more accurately; it does not by
itself recover the real-line operator.  Boundary truncation, interior discretization,
retained spectral dimension, coordinate injection, eigensolver error, and floating-point
error therefore remain distinct \cite{strikwerda2004}.  This is the same logical reason
that agreement inside a reduced model cannot repair an inadequate parent.

For the silicon hierarchy, four notions of ``finite'' must consequently not be merged:
\begin{enumerate}[leftmargin=2em]
 \item a finite plane-wave cutoff and reciprocal sampling define a numerical DFT
 representation, even though each plane wave extends across the cell;
 \item a finite band/Wannier rank defines a retained spectral state space;
 \item a finite real-space hopping support defines the range of the represented lattice
 Hamiltonian; and
 \item a finite evaluation inventory defines only the reciprocal points on which a loss
 has been sampled.
\end{enumerate}
Taking the plane-wave cutoff, construction mesh, retained rank, hopping radius, or
evaluation mesh to a larger value addresses a different approximation.  In particular,
restoring all native hopping blocks converges to the fixed Wannier source, not to its DFT
parent; refining a support-limited model cannot remove a source/parent discrepancy.
This nested-limit view is the conceptual bridge from the PIAB1D and QHO1D calibration
reports to the present periodic reduction study.

\subsection{Gauge freedom and locality}

Wannier functions inherit the gauge freedom of a retained Bloch subspace
\cite{marzari1997,souza2001,marzari2012}.  Smooth gauges can differ in spread and real-space
shape even when they span the same projector.  For a Fourier-expanded transformation
$U(\kvec)=\sum_{\mathbf m}e^{i\kvec\cdot\mathbf m}U_{\mathbf m}$, substitution into
Eq.~\eqref{eq:intro-gauge} gives the hopping convolution
\begin{equation}
 \widetilde T_{\mathbf Q}
 =\sum_{\mathbf m,\mathbf n}
 U_{\mathbf m}^\dagger T_{\mathbf Q+\mathbf m-\mathbf n}U_{\mathbf n}.
 \label{eq:hopping-convolution}
\end{equation}
A constant transformation has only one Fourier coefficient and does not enlarge support.
A momentum-dependent transformation mixes displacement indices.  Locality is therefore
a property of the represented field and allowed gauge, not of the complete spectrum
alone.

This does not make locality arbitrary.  Once the construction, retained space, gauge,
metric, support family, and error criterion are fixed, the resulting locality problem is
well defined.  The important restriction is epistemic: spectral data alone cannot
identify a unique hopping distribution.

\subsection{Gauge invariants, alignment families, and identifiability}

For each momentum, unitary similarity preserves the characteristic polynomial,
trace invariants, and the complete unordered spectrum:
\begin{equation}
 \det[zI-U(\kvec)^\dagger H(\kvec)U(\kvec)]
 =\det[zI-H(\kvec)].
 \label{eq:characteristic-invariance}
\end{equation}
It also preserves every pointwise unitarily invariant norm of the complete operator.
It does not preserve individual matrix entries, basis-resolved observables, or Fourier
coefficients.  Consequently, an eigenvalue-only inverse problem has an entire gauge
orbit of solutions.  A locality statement becomes identifiable only after the allowed
coordinate transformations are restricted.

The present work distinguishes four such restrictions.  Identity alignment compares
matrices in the declared common Wannier basis.  Global alignment admits one momentum-independent unitary
$V\in U(N)$ and minimizes a weighted aggregate residual,
\begin{equation}
 \mathcal L_{\rm global}(V)^2
 =\sum_j w_j\norm{H(\kvec_j)-V^\dagger A(\kvec_j)V}_F^2.
 \label{eq:global-alignment}
\end{equation}
Pointwise polar alignment admits a different complete unitary at every point.  If two
orthonormal frame matrices $X$ and $Y$ have overlap $X^\dagger Y=L\Sigma R^\dagger$,
the polar/Procrustes action is $Q=LR^\dagger$ when the declared singular-value rank gate
is satisfied \cite{schonemann1966,higham1986}.  Finally, a known-inverse action uses the
unitary supplied by a synthetic construction.  The latter two are oracle controls: they
check the representation and implementation, but they are not available from the
complete spectra alone.

The minimized value in Eq.~\eqref{eq:global-alignment} is a distance under this
particular weighted global-unitary action; it is not asserted to be an invariant metric
between unrestricted gauge orbits.  A nonzero identity or global residual is therefore
not evidence that two complete unitarily related fields represent different physics.
It states that the generating momentum dependence lies outside the admitted alignment
family.  Conversely, allowing
an independently fitted pointwise unitary in a same-gauge truncation study would erase
part of the approximation being measured.  This is why the shell-reduction operator
loss uses identity alignment.

\subsection{Norm relations, Parseval identity, and support truncation}

For Hermitian matrices $A$ and $B$, the ordered eigenvalues satisfy the
Hoffman--Wielandt bound
\begin{equation}
 \sum_{n=1}^{N}[\lambda_n(A)-\lambda_n(B)]^2
 \leq \norm{A-B}_F^2.
 \label{eq:hoffman-wielandt}
\end{equation}
Thus an operator residual controls an ordered spectral residual at the same momentum.
The converse is false: isospectral matrices can have a finite Frobenius separation.
This one-way implication motivates retaining both loss coordinates rather than treating
one as a surrogate for the other.

For two matrix fields in the same Fourier convention,
$\Delta H(\kvec)=\sum_{\Rvec}e^{i\kvec\cdot\Rvec}\Delta T(\Rvec)$, normalized
Brillouin-zone integration gives the matrix-valued Parseval identity
\begin{equation}
 \frac{1}{\Omega_{\rm BZ}}\int_{\rm BZ}
 \norm{\Delta H(\kvec)}_F^2\,d\kvec
 =\sum_{\Rvec}\norm{\Delta T(\Rvec)}_F^2.
 \label{eq:matrix-parseval}
\end{equation}
For rank $N$ and reporting energy $E_0$, the exact integrated per-state operator RMS is
\begin{equation}
 \tau_{2,{\rm BZ}}
 =\frac{1}{E_0}\left[
 \frac{1}{N}\sum_{\Rvec}\norm{\Delta T(\Rvec)}_F^2
 \right]^{1/2}.
 \label{eq:parseval-rms}
\end{equation}
No reciprocal sampling is required once the unique coefficient convention is fixed.
This exact identity is used in the fidelity ladder.

For a literal support $S$, the truncated model is
\begin{equation}
 H_S(\kvec)=\sum_{\Rvec\in S}e^{i\kvec\cdot\Rvec}T(\Rvec),
 \qquad
 \Delta T_S(\Rvec)=
 \begin{cases}
 0,&\Rvec\in S,\\
 -T(\Rvec),&\Rvec\notin S.
 \end{cases}
 \label{eq:literal-truncation}
\end{equation}
Equation~\eqref{eq:parseval-rms} then makes the integrated operator loss exactly the
excluded Frobenius coefficient power.  A separate uniform operator bound follows from
the triangle inequality,
\begin{equation}
 \sup_{\kvec}\norm{H_S(\kvec)-H(\kvec)}_2
 \leq \sum_{\Rvec\notin S}\norm{T(\Rvec)}_2.
 \label{eq:triangle-tail-bound}
\end{equation}
The Frobenius tail and the spectral-norm triangle bound answer different questions: the
first is an exact integrated RMS quantity, while the second is generally a conservative
uniform bound.  Neither determines the local curvature of an individual band.

\subsection{Spectral and operator losses}

For a reduced model $H_{\rm r}(\kvec)$ in the same basis as a reference
$H_{\rm W}(\kvec)$, a spectral loss compares ordered eigenvalues, while an operator loss
compares matrices.  On the silicon 408-point inventory the primary normalized RMS
coordinates are
\begin{align}
 \tau_1^2
 &=\frac{1}{4080}\sum_{\kvec,n}
 \left(\frac{\epsilon^{\rm r}_{n\kvec}-\epsilon^{\rm W}_{n\kvec}}
 {1\,{\rm eV}}\right)^2,\label{eq:tau1-main}\\
 \tau_2^2
 &=\frac{1}{4080}\sum_{\kvec}
 \frac{\norm{H_{\rm r}(\kvec)-H_{\rm W}(\kvec)}_F^2}
 {(1\,{\rm eV})^2}.
 \label{eq:tau2-main}
\end{align}
The 1 eV denominator is a reporting normalization, not an acceptance limit.  Identity
alignment is used for same-gauge shell models because a fitted pointwise alignment would
erase part of the reduction error.

A matrix perturbation bounds spectral displacement, but spectral agreement does not
recover the matrix.  Consequently, $(\tau_1,\tau_2)$ is treated as a two-dimensional
loss coordinate.  Post-hoc threshold maps summarize observed witnesses; they do not
become prospective scientific tolerances after the fact.

\subsection{Global averages and band-edge observables}

A global RMS loss averages across bands and reciprocal points.  A band edge depends on a
local extremum and its derivatives.  For a nondegenerate band,
\begin{equation}
 (m^{-1})_{ij}=\frac{1}{\hbar^2}
 \frac{\partial^2E}{\partial k_i\partial k_j}.
 \label{eq:mass-tensor}
\end{equation}
Differentiation can amplify small local shape errors that contribute little to a global
average.  Valley displacement, Hessian rotation, and principal masses are therefore
first-class reduction outputs rather than decorations on a band plot.

At a degenerate valence edge, independently ordered branches may exchange identity and
fail to be analytic.  A single ``hole mass tensor'' is then not basis invariant.  The
appropriate object is a coupled matrix-valued quadratic Hamiltonian, consistent with
multiband effective-mass reasoning \cite{LuttingerKohn1955}.  In the present non-SOC
source, this object describes three orbital valence channels; it is not assigned physical
heavy-hole, light-hole, or split-off labels.

Let $P_0$ project onto an $r$-dimensional cluster at $\kvec_0$.  A smooth nearby cluster
with projector $P(\qvec)$ must first be identified with $P_0\C^N$ before matrices can be
subtracted.  If $X_0$ and $X(\qvec)$ are orthonormal cluster frames, the singular values
of $X_0^\dagger X(\qvec)$ are cosines of principal angles.  The polar factor of this
overlap supplies a least-squares unitary identification whenever the overlap remains
full rank.  The transported represented cluster Hamiltonian can then be expanded as
\begin{equation}
 H_{\rm eff}(\qvec)=H_{\rm eff}(0)
 +\sum_i q_i A_i+\sum_{i,j}q_iq_jB_{ij}+O(\norm{\qvec}^3).
 \label{eq:general-cluster-expansion}
\end{equation}
A symmetric $+\qvec/-\qvec$ average removes the odd term without assigning individual
branch identities.  For a unit direction $\hat{\mathbf n}$, the eigenvalues of
$2\sum_{ij}n_in_jB_{ij}$ are the quadratic curvature channels.  These channels are
basis invariant under one common unitary rotation of the cluster coefficient field.
They need not equal Hessians of independently ordered bands when $H_{\rm eff}(0)$ is
degenerate or when a reduction has artificially split it.

\section{Exact rank-two theory}

\subsection{Trace--traceless candidate family}

Write a rank-two represented Hamiltonian as
\begin{equation}
 H(q)=\bar E(q)I_2+D(q),\qquad \operatorname{tr}D(q)=0.
 \label{eq:trace-traceless}
\end{equation}
The candidate family shifts the centroid and positively rescales the splitting,
\begin{equation}
 C_{s,\lambda}(q)=\bar E(q)I_2+sE_{\rm ref}I_2+\lambda D(q),
 \qquad \lambda>0.
 \label{eq:candidate-rank2}
\end{equation}
The comparison target is a pointwise band-equivalent field
\begin{equation}
 A_a(q)=U_a(q)^\dagger H(q)U_a(q),
 \qquad U_a(q)=R(a\sin2\pi q).
 \label{eq:smooth-attack}
\end{equation}
The admissible alignment is one momentum-independent real rotation $R(\phi)$, not the
generating momentum-dependent map.

With positive normalized weights $\alpha_i$, define
\begin{equation}
 \kappa=\sum_i\frac{\alpha_i}{2E_{\rm ref}^2}\norm{D(q_i)}_F^2.
 \label{eq:kappa-main}
\end{equation}
The losses use the same positive weights, normalized by $\sum_i\alpha_i=1$:
\begin{align}
 L_{\rm spec}^2&=\sum_i\frac{\alpha_i}{2E_{\rm ref}^2}
 \sum_{n=1}^{2}[\epsilon_n(C_{s,\lambda};q_i)-\epsilon_n(A_a;q_i)]^2,\\
 L_{{\rm op},\phi}^2&=\sum_i\frac{\alpha_i}{2E_{\rm ref}^2}
 \norm{R(\phi)^\mathsf{T}C_{s,\lambda}(q_i)R(\phi)-A_a(q_i)}_F^2.
 \label{eq:rank2-loss-definitions}
\end{align}
For an allowed global angle, let $\mu_\phi$ be the normalized weighted correlation
between the globally rotated candidate traceless field and the attacked target.  A real
rotation conjugates the two-dimensional traceless symmetric coordinates through angle
$2\phi$.  Consequently the correlation has the explicit form
\begin{equation}
 \kappa\mu_\phi=b_0+c\cos2\phi+d\sin2\phi,
 \qquad
 \kappa\mu_{\max}=b_0+\sqrt{c^2+d^2}.
 \label{eq:second-harmonic-correlation}
\end{equation}
This is why the global alignment reduces to one second harmonic rather than a numerical
search over an unrestricted field of rotations.

\begin{theorem}[Common-curvature loss geometry]
Under the family above,
\begin{align}
 L_{\rm spec}^2(s,\lambda)
 &=s^2+\kappa(\lambda-1)^2,\label{eq:spec-quadratic}\\
 L_{\rm op,*}^2(s,\lambda)
 &=s^2+\kappa(\lambda-\mu_{\max})^2
 +\kappa(1-\mu_{\max}^2).
 \label{eq:op-quadratic}
\end{align}
The two losses have common curvature $\operatorname{diag}(1,\kappa)$.  Their full
unclipped admissible sets are homothetic ellipses with different centers and, for the
operator channel, an irreducible floor $\kappa(1-\mu_{\max}^2)$.
\end{theorem}

The theorem makes the compatibility question explicit.  Spectral adequacy centers the
splitting scale at $\lambda=1$; restricted operator adequacy centers it at
$\lambda=\mu_{\max}$.  The momentum-dependent part that cannot be represented by one
global rotation remains as the irreducible floor.

If the spectral and operator thresholds are $\tau_s$ and $\tau_o$, define
\begin{equation}
 r_*^2=\tau_o^2-\kappa(1-\mu_{\max}^2).
\end{equation}
When the complete threshold ellipses are contained in the declared parameter rectangle,
they intersect if and only if
\begin{equation}
 \sqrt\kappa\,|1-\mu_{\max}|\le \tau_s+r_*.
 \label{eq:ellipse-intersection}
\end{equation}
The containment condition is necessary because
Eq.~\eqref{eq:ellipse-intersection} compares the full ellipses through their center
separation and radii.  If the admissible rectangle clips either ellipse, two full
ellipses may intersect only outside the allowed parameter domain, or clipping may remove
the nearest points used by the center-distance argument.  A failed numerical search is
therefore not a separation proof; Eq.~\eqref{eq:ellipse-intersection} supplies the
analytic decision only together with the containment premises.

\subsection{Smooth threshold continuum}

For a flat real-symmetric traceless parent, the smooth attack in
Eq.~\eqref{eq:smooth-attack} yields
\begin{equation}
 \mu_{\max}(a)=\frac{1}{2\pi}\int_0^{2\pi}
 \cos[2a\sin\theta]\,d\theta=J_0(2a),
 \label{eq:bessel-correlation}
\end{equation}
where $J_0$ is the Bessel function of the first kind \cite{dlmf1012}.  The factor two
appears because conjugation by a real rotation through $\theta_a$ rotates the traceless
symmetric coordinates through $2\theta_a$.  Therefore
\begin{align}
 \tau_{o,\min}(a)
 &=\sqrt{\kappa[1-J_0(2a)^2]},\label{eq:min-threshold}\\
 \tau_{o,\rm joint}(a;0)
 &=\sqrt{2\kappa[1-J_0(2a)]}.
 \label{eq:joint-threshold}
\end{align}
On $0\le a\le\pi/4$, the Bessel argument satisfies $0\le2a\le\pi/2$, before the
first positive zero of $J_1$.  Since $J_0'(x)=-J_1(x)<0$ on this interval,
$J_0(2a)$ decreases and both thresholds are nondecreasing.  Every member of the family
has exactly the same complete bands; only the restricted representation threshold
changes.

\subsection{Residual--locality identity}

Let the flat traceless target after attack be
\begin{equation}
 B(q)=U(q)^\dagger D U(q)=\sum_R e^{2\pi iqR}B_R.
\end{equation}

\begin{theorem}[Rank-two residual--non-onsite-power identity]
For the real rotation orbit used above,
\begin{equation}
 \inf_{s,\lambda\ge0,\phi}L_{\rm op,\phi}^2
 =\kappa(1-\mu_{\max}^2)
 =\frac{1}{2E_{\rm ref}^2}\sum_{R\ne0}\norm{B_R}_F^2.
 \label{eq:residual-tail-identity}
\end{equation}
The same identity holds for the declared finite uniform transform.
\end{theorem}

To expose the intermediate step, normalized Parseval gives
\begin{equation}
 \sum_{R\ne0}\norm{B_R}_F^2
 =\int_{-1/2}^{1/2}\norm{B(q)}_F^2\,dq-\norm{B_0}_F^2.
 \label{eq:rank2-parseval-intermediate}
\end{equation}
The best constant real rotation and positive scale reproduce the optimally oriented
onsite component; the normalized squared onsite amplitude is $\mu_{\max}^2$.  Because
unitary conjugation preserves $\norm{B(q)}_F=\norm D_F$, the remainder is exactly
$\norm D_F^2(1-\mu_{\max}^2)$, yielding
Eq.~\eqref{eq:residual-tail-identity}.  This is the exact bridge between restricted
matrix mismatch and hopping locality under the theorem's assumptions.  The global
rotation can optimize the onsite orientation, but it cannot remove Fourier weight at
nonzero displacement.  For the smooth family, the threshold in
Eq.~\eqref{eq:min-threshold} is the square root of the normalized induced non-onsite
power.

\subsection{What the theorem does not claim}

The result does not establish a universal relation between any matrix residual and any
locality measure.  It assumes a flat traceless rank-two parent and matched real rotation
orbits.  A dispersive parent adds native nonzero-range blocks and cross terms.  Higher
rank introduces noncommuting rotations and a richer gauge group.  The value of the
controlled theorem is that it isolates a mechanism that can then be tested without
silently extending the equality beyond its assumptions.

\section{Controlled rank-two calculation}

\subsection{Finite-range four-state parent}

The dispersive benchmark is a one-dimensional periodic model with lattice period $a_0$
and reduced momentum $q=ka_0/(2\pi)\in[-1/2,1/2)$.  Four states per cell are partitioned
into lower and upper two-state sectors.  The only nonzero parent hopping blocks are
\begin{align}
 T_0&=\begin{pmatrix}D_{\rm L}&0\\0&D_{\rm U}\end{pmatrix},&
 T_{+1}&=\begin{pmatrix}J_{\rm L}&V\\0&J_{\rm U}\end{pmatrix},&
 T_{-1}&=T_{+1}^\dagger,
 \label{eq:rank2-parent-blocks}
\end{align}
with, in units of $E_{\rm ref}$,
\begin{align}
 D_{\rm L}&=\begin{pmatrix}-1.20&0.06\\0.06&-0.35\end{pmatrix},&
 D_{\rm U}&=\begin{pmatrix}1.00&0.04\\0.04&1.75\end{pmatrix},\nonumber\\
 J_{\rm L}&=\begin{pmatrix}0.15&0\\0&-0.10\end{pmatrix},&
 J_{\rm U}&=\begin{pmatrix}0.05&0\\0&-0.04\end{pmatrix},&
 V&=\begin{pmatrix}0.12&0\\0&0.10\end{pmatrix}.
 \label{eq:rank2-parent-values}
\end{align}
The Bloch matrix is
\begin{equation}
 K(q)=\begin{pmatrix}
 D_{\rm L}+zJ_{\rm L}+z^{-1}J_{\rm L}^\dagger&zV\\
 z^{-1}V^\dagger&D_{\rm U}+zJ_{\rm U}+z^{-1}J_{\rm U}^\dagger
 \end{pmatrix},\qquad z=e^{2\pi iq}.
 \label{eq:rank2-parent-bloch}
\end{equation}
The parent has exact range one.  Its intersector coupling makes the lower two-dimensional
spectral projector momentum dependent even though the parent itself is finite range.
$E_{\rm ref}$ is only a normalization scale.

At every sampled momentum, the four Hermitian eigenpairs are sorted and the two lowest
states are retained.  The second-to-third band gap is monitored so that the retained
rank and band identity are not changed silently.  Raw eigensolver vectors are not used
as a common coordinate system because their phases and internal rank-two rotations are
independent at each point.

\subsection{Parallel-transport frame and periodic sewing}

Let $Y_i\in\C^{4\times2}$ contain the raw retained eigenvectors at $q_i$, and let $F_i$
be the accepted transported frame.  Starting with $F_0=Y_0$, form
\begin{equation}
 M_i=F_i^\dagger Y_{i+1}=L_i\Sigma_iR_i^\dagger,
 \qquad F_{i+1}=Y_{i+1}R_iL_i^\dagger.
 \label{eq:rank2-parallel-transport}
\end{equation}
The rotation minimizes $\norm{Y_{i+1}Q-F_i}_F$ over unitary $Q$.  Both singular values
of $M_i$ must pass the declared rank gate; a rank-deficient overlap is rejected instead
of being completed arbitrarily.  Orthonormality and retained-projector reconstruction
are checked at every node.

Transport around the periodic mesh generally leaves a sewing mismatch.  The duplicate
endpoint is transported once more, the unitary polar part of
$S=F_0^\dagger F_N$ is diagonalized as
$S=Z\operatorname{diag}(e^{i\gamma_1},e^{i\gamma_2})Z^\dagger$, and the mismatch is
distributed through
\begin{equation}
 W_i=Z\operatorname{diag}(e^{-i(i/N)\gamma_1},e^{-i(i/N)\gamma_2})Z^\dagger,
 \qquad F_i\leftarrow F_iW_i.
 \label{eq:rank2-sewing}
\end{equation}
The corrected endpoint must equal the start within tolerance.  The retained represented
Hamiltonian is
\begin{equation}
 H(q_i)=F_i^\dagger K(q_i)F_i.
 \label{eq:rank2-retained-field}
\end{equation}
The attack $U_a(q)=R[a\sin(2\pi q)]$ then gives
$A_a(q)=U_a(q)^\dagger H(q)U_a(q)$ for $0\le a\le\pi/4$.  Complete spectra are equal by
construction, while the comparison admits only one subsequent global real rotation.

\subsection{Finite Fourier representation and range decision}

For an $N$-point centered construction mesh,
\begin{equation}
 T_R^{(N)}=\frac{1}{N}\sum_{i=0}^{N-1}e^{-2\pi iq_iR}H(q_i),
 \qquad H_N(q)=\sum_R e^{2\pi iqR}T_R^{(N)}.
 \label{eq:rank2-finite-fourier}
\end{equation}
Construction-node reconstruction, cyclic Hermiticity, and Parseval equality are checked.
For even $N$, the centered Nyquist representative $-N/2$ is self-conjugate only in the
cyclic convention.  It is removed before ordinary-lattice truncation, and its power is
reported as a conversion channel.  The ordinary inventory is
$\{-N/2+1,\ldots,N/2-1\}$ and must satisfy $T_{-R}=T_R^\dagger$.

For each integer cutoff $R_c$, the model
\begin{equation}
 H_{N,R_c}^{\rm ord}(q)=\sum_{|R|\le R_c}e^{2\pi iqR}T_R^{(N)}
 \label{eq:rank2-truncated-field}
\end{equation}
is evaluated on a disjoint reciprocal inventory.  Matrix maximum/RMS and ordered-band
maximum/RMS errors are retained separately for interpolation, deliberate truncation,
and total finite-model consequences.  The first-passing range is the smallest tested
integer $R_c$ whose sampled sorted-band error does not exceed a declared target.  The
search is exact only on its bounded inventory; it is not a continuum localization
length.

The flat analytic family is evaluated directly from the Bessel expressions.  The
non-flat parent uses directly sampled attack amplitudes rather than an interpolated
continuum claim: 65 amplitudes for momentum-refinement diagnostics and 33 amplitudes for
the parent/truncation surfaces.  Momentum resolutions from 32 through 512 are retained,
and the headline endpoint uses $N=512$ with cutoffs $R_c=0,\ldots,16$.

\subsection{Synthetic numerical results}

At the flat endpoint $a=\pi/4$, the normalized minimum operator and exact-spectrum joint
thresholds are approximately 0.8815978972 and 1.0276174232.  The induced non-onsite
power is approximately 0.7772149, consistent with the squared minimum threshold.  The
finite transform is roundoff-limited over the tested $N=32,\ldots,512$ sequence.

For the dispersive parent at $N=512$, the sampled external gap is approximately
$1.0566461E_{\rm ref}$ and the minimum transport-overlap singular value is
0.9999997323.  Complete bands agree to about $1.45\times10^{-15}E_{\rm ref}$, while the
comparison non-onsite power rises from 0.0323839143 to 0.1761532501 over the attack.

The central finite-range result is measured at equal sampled accuracy.  At a band-error
target of $10^{-8}E_{\rm ref}$, the first passing cutoffs are
\begin{equation}
 \widehat R_{c,\rm ref}^{\star}=4,
 \qquad
 \widehat R_{c,\rm comp}^{\star}=10.
\end{equation}
The band-equivalent comparison therefore requires six additional lattice steps on the
declared cutoff inventory.  At targets $10^{-6}$, $10^{-8}$, $10^{-10}$, and
$10^{-12}E_{\rm ref}$, the endpoint reference/comparison pairs are
$(3,8)$, $(4,10)$, $(6,12)$, and $(7,14)$.

\resultfigure
 {../../../../../calculations/periodic1d/rank2/band-equivalence-versus-hopping-locality/analysis/development/figures/band-equivalence-locality-scientific-story.pdf}
 {Controlled rank-two mechanism from exact band equivalence to representation-dependent
 hopping locality.}
 {The complete fields remain pointwise unitary equivalents; locality differences become
 finite-range consequences only after a support rule and accuracy criterion are imposed.}
 {fig:rank2-story}

\resultfigure
 {../../../../../calculations/periodic1d/rank2/band-equivalence-versus-hopping-locality/analysis/development/figures/range-inflation-preview.pdf}
 {Sampled fixed-accuracy range inflation across parent scale, attack amplitude, and
 error target.}
 {The first-passing search is exact only over the declared integer cap
 $R_c=0,\ldots,16$; failure within the cap is not a universal no-solution result.}
 {fig:range-inflation}

\resultfigure
 {../../../../../calculations/periodic1d/rank2/band-equivalence-versus-hopping-locality/analysis/development/figures/truncation-error-channel-decomposition.pdf}
 {Separated finite-representation error channels in the rank-two calculation.}
 {Cyclic conversion, interpolation, deliberate truncation, and total consequences are
 displayed independently and are not treated as additive spectral errors.}
 {fig:rank2-error-channels}

\resultfigure
 {../../../../../calculations/periodic1d/rank2/band-equivalence-versus-hopping-locality/analysis/development/figures/frame-sewing-diagnostics.pdf}
 {Retained-frame transport and sewing diagnostics for the four-state parent.}
 {The external gap, neighboring-frame overlap, orthonormality, projector reconstruction,
 and periodic closure are gates on the represented rank-two field.}
 {fig:rank2-sewing}

\resultfigure
 {../../../../../calculations/periodic1d/rank2/band-equivalence-versus-hopping-locality/analysis/development/figures/hamiltonian-spectral-diagnostics.pdf}
 {Hamiltonian and complete-spectrum diagnostics in the controlled calculation.}
 {The figure distinguishes matrix-field change from eigenvalue invariance before any
 finite-range approximation is imposed.}
 {fig:rank2-hamiltonian-diagnostics}

\resultfigure
 {../../../../../calculations/periodic1d/rank2/band-equivalence-versus-hopping-locality/analysis/development/figures/sampled-range-decision-map.pdf}
 {Sampled first-passing range decisions over the declared attack and parent-scale grid.}
 {Each cell reports a bounded integer search result, not a continuum optimum or a
 universal localization length.}
 {fig:rank2-range-map}

\resultfigure
 {../../../../../calculations/periodic1d/rank2/band-equivalence-versus-hopping-locality/analysis/development/figures/truncation-band-consequences.pdf}
 {Band consequences of finite-range truncation in the controlled family.}
 {Reference and attacked representations are assessed at identical support rules and
 evaluation inventories.}
 {fig:rank2-truncation-bands}

The active rank-two calculation has a passing unsealed direct-v2 execution with
byte-owning independent verification for its declared final-mesh parent family,
scale-one truncation slice, and refinement channels.  It does not authenticate every
five-scale development surface used in the broad preview figures.  The exact theorem is
not weakened by that lifecycle limitation; the numerical publication headline remains
development or partial evidence until a complete-scope retained package is separately
authorized and verified.

\section{Silicon source and reduction methodology}

\subsection{Source construction and adequacy boundary}

The material-facing source is the authenticated $C6$-B20-W21 rank-10 Wannier Hamiltonian
constructed from scalar-relativistic, non-spin-polarized PBE silicon calculations using
Quantum ESPRESSO 7.5 and Wannier90 3.1.0 \cite{giannozzi2009,pizzi2020}.  The notation
records a Gamma-origin $6^3$ construction mesh, 20 parent bands, and a rank-10 retained
space.  Spinors are absent.  The analysis therefore concerns a scalar-relativistic
orbital Hamiltonian, not an SOC spinor Hamiltonian.

The primitive cell contains Si atoms at fractional coordinates $(0,0,0)$ and
$(1/4,1/4,1/4)$.  The relaxed right-handed FCC primitive basis, in angstrom, is
\begin{equation}
 A=\begin{pmatrix}
 2.73459104682&0&2.73459104682\\
 2.73459104682&2.73459104682&0\\
 0&2.73459104682&2.73459104682
 \end{pmatrix},
 \label{eq:silicon-direct-basis}
\end{equation}
where the primitive vectors are the columns of $A$, and the corresponding conventional
cubic lattice parameter is
$a=5.469182093640207$~\AA.  The reciprocal-basis columns are
\begin{equation}
 B=\begin{pmatrix}
 1.14883454228&-1.14883454228&1.14883454228\\
 1.14883454228&1.14883454228&-1.14883454228\\
 -1.14883454228&1.14883454228&1.14883454228
 \end{pmatrix}\ \text{\AA}^{-1}.
 \label{eq:silicon-reciprocal-basis}
\end{equation}
The cell was obtained from an authorized FCC-native \texttt{vc-relax} with
\texttt{ibrav=2}, initial $A=5.43$~\AA, and scalar-cell relaxation.  Three BFGS steps
and four structural SCF cycles reached a reported terminal pressure of 0.01~kbar and a
final-SCF pressure of 0.09~kbar under a 0.5~kbar declared tolerance.  These observations
establish the retained geometry occurrence; they do not constitute a general structural
uncertainty analysis.

The accepted parent SCF used the SSSP 1.3.0 PBE-efficiency silicon ultrasoft
pseudopotential, a Gamma-centered $10^3$ mesh, and wavefunction/density cutoffs of
42/336~Ry.  The retained occurrence reported electronic convergence in five iterations.
The Wannier parent NSCF used the same accepted structure and cutoffs, 216 ordered
Gamma-origin $6^3$ points, 20 bands, and Davidson diagonalization.  Its eigenvalue range
was $[-5.83413,30.27201]$~eV.  These calculations were performed and authorized in their
own retained campaign records; they were not rerun for the present post-hoc study.

Wannier90 retained ten functions from the 20-band inventory.  The input used
\texttt{spinors=.false.}, outer disentanglement window $[-6,21]$~eV, frozen window
$[-5.85,8]$~eV, and projections \texttt{Si:sp3} together with one additional radial
$s$-like projection per Si specified by \texttt{Si:l=0:r=2}.  Both disentanglement and
Wannier localization had 40000-step ceilings and $2\times10^{-9}$ convergence gates.
The retained run converged at disentanglement iteration 1517 and Wannier iteration 378,
with final total spread $21.022751219$~\AA$^2$.  This is native optimizer convergence,
not evidence of a global spread minimum.

The native HR file contains 279 Wannier90 representatives.  After applying the retained
Wigner--Seitz-vector/degeneracy convention, the effective ordinary reconstruction
contains 383 representatives.  The represented field is
\begin{equation}
 H_{\rm W}(\kvec)=\sum_{\Rvec\in\mathcal R_{\rm WS}}
 d_{\Rvec}^{-1}e^{i\kvec\cdot\Rvec}H_{\rm HR}(\Rvec),
 \label{eq:native-wannier-reconstruction}
\end{equation}
where the effective coefficient inventory incorporates the authenticated Wigner--Seitz
mapping and degeneracy factors $d_{\Rvec}$.  Pair closure, Hermiticity, construction-mesh
reconstruction, reciprocal periodicity, and finite spectra were checked before the
source was admitted to reduction analysis.

Structural eligibility does not imply material adequacy.  In the canonical corrected
current-source comparison against the revealed 408-point DFT parent, the source indirect
gap is 0.6892447~eV and the parent gap is 0.6112467~eV, a 12.7605\% relative difference
against the prospectively declared 5\% target.  All six valley positions pass their
0.034465~\AA$^{-1}$ displacement target, with displacements from approximately 0.01906
to 0.02134~\AA$^{-1}$, but all six longitudinal masses and eleven of twelve transverse
masses fail the 8\% target.  The non-gating injective 408-point spectral residual is
0.246196~eV RMS and 0.737286~eV maximum.  These failures bound every result below:
source-relative preservation is not DFT validation.

An earlier parent-comparison record reported a 0.5226105~eV candidate gap.  The retained
source identities resolve the apparent discrepancy: that value belongs to a distinct
lower-energy-frozen $4^3$/12-band Wannier candidate whose eligible HR hash is
\texttt{2139eb7f...}, not to the present $C6$--$B20$--$W21$ candidate whose HR hash is
\texttt{24a411c8...}.  The two records share the same revealed DFT parent gap but do not
share a candidate Hamiltonian.  The present article therefore uses 0.6892447~eV only for
the authenticated $C6$--$B20$--$W21$ source and does not combine the candidate lineages.

No frozen 402-point confirmation values are used.  No Quantum ESPRESSO, interface, or
Wannier90 process was invoked for the gauge, shell, fidelity, or band-edge analyses.
Every subsequent Hamiltonian is reconstructed from the same authenticated HR and
Wigner--Seitz data in the same rank-10 Wannier basis and native energy reference.

For band-structure context, one separately authorized fixed-density QE 7.5 calculation
was performed after the referee revision.  It consumed an isolated copy of the accepted
K10 SCF state and evaluated 20 parent Kohn--Sham bands at 241 ordered points on
$L$--$\Gamma$--$X$--$\Gamma$.  Figure~\ref{fig:parent-dft-bands} displays the lowest 12
parent bands together with the complete rank-10 Wannier source.  Both curves retain their
native common energy reference; the plot subtracts the parent path VBM and applies no
fitted source shift.  The sampled path gaps are 0.6114~eV for the parent and
0.69046~eV for the source.  These rounded path diagnostics visualize, but do not replace,
the separately optimized indirect-gap records.

\resultfigure
 {../../../../../calculations/Si/bulk/simulations/evidence/sssp-1.3.0-pbe-efficiency-qe-7.5-accepted-s1-parent-bands-l-gamma-x-gamma-b20/figures/parent-dft-band-diagram.pdf}
 {Accepted-S1 parent DFT band path and full Wannier source.}
 {The solid curves are the lowest 12 of 20 scalar-relativistic PBE parent bands; dashed
 curves are the ten bands of the authenticated $C6$--$B20$--$W21$ source.  The horizontal
 path contains 241 ordered points.  The common display zero is the parent path VBM, and
 no fitted energy alignment is used.}
 {fig:parent-dft-bands}

\subsection{Higher-rank gauge attack}

The internal attack is frozen independently of the silicon Hamiltonian.  It is a real
orthogonal---and therefore complex-unitary---rank-10, depth-2 alternating
nearest-neighbor Givens circuit in the authenticated Wannier coordinate order.  The
indices below refer only to that retained coordinate order; no atomic-orbital identity
is assigned to an index.
\begin{equation}
 G(\kvec)=G_1(\kvec)G_2(\kvec),\qquad
 A(\kvec)=G(\kvec)^\dagger H_{\rm W}(\kvec)G(\kvec).
 \label{eq:rank10-attack}
\end{equation}
Layer 1 rotates pairs $(0,1),(2,3),(4,5),(6,7),(8,9)$; layer 2 rotates
$(1,2),(3,4),(5,6),(7,8)$.  Let the harmonic bonds be
$\mathbf b_1=(1,0,0)$, $\mathbf b_2=(0,1,0)$, and
$\mathbf b_3=(0,0,1)$.  The shared layer scalar is
\begin{align}
 g_\ell(\kvec)=\frac{1}{3}\sum_{a=1}^{3}
 \left[\frac{3}{4}\sin(2\pi\kvec\cdot\mathbf b_a)
 +\frac{1}{4}(-1)^\ell\sin(4\pi\kvec\cdot\mathbf b_a)\right].
 \label{eq:gauge-scalar}
\end{align}
For a pair whose right index is $p$, the rotation angle is
\begin{equation}
 \theta_{\ell p}(\kvec)=\frac{\pi}{8}
 (-1)^{\ell+p}\frac{1+0.15(2p-10)/9}{1.15}\,g_\ell(\kvec).
 \label{eq:gauge-angle}
\end{equation}
This construction is periodic, has conservative angle bound $\pi/8$, and supplies its
exact inverse $G(\kvec)^\dagger$.  It was selected as a synthetic middle-amplitude
control before inspecting the silicon support results; changing a bond, depth,
normalization, or coefficient defines a different attack.  It is a constructed stress
test demonstrating one allowed redistribution of locality.  It is not claimed to be
representative of typical Wannier90 gauge variation or to maximize range inflation.

The source construction coordinates are the native Gamma-origin points
$0,1/6,\ldots,5/6$ on each reduced reciprocal axis.  For the centered transform they
are wrapped and permuted to
\begin{equation}
 \mathcal K_6=\{-1/2,-1/3,-1/6,0,1/6,1/3\}^{\otimes3},
 \label{eq:centered-c6-grid}
\end{equation}
with zero construction offset.  The direct native HR/Wigner--Seitz route is evaluated
first; the generic field adapter does not reconstruct the source from raw DFT
wavefunctions.  The distinct evaluation inventory is the ordered 408-point development
geometry with uniform weights $w_j=1/408$.  It is accepted as an arbitrary point list
and is not padded into a tensor grid.  The frozen 402-point confirmation geometry and
values play no role.

Four alignment channels are evaluated.  Identity uses the common reference coordinates.
Global-unitary alignment fits one complete momentum-independent unitary over the
construction field.  Pointwise polar alignment obtains one unitary per coordinate from
the admitted frame overlaps.  Known-inverse alignment applies the exact inverse action
of Eq.~\eqref{eq:rank10-attack}.  The frozen rank cutoffs use the additive rule with
absolute threshold $10^{-12}$ and relative threshold $10^{-10}$; the same absolute and
relative values are used for transport-rank checks.  A deficient overlap fails rather
than receiving an implicit null-space completion.  Pointwise polar and known inverse
are oracle controls, not candidate discovery procedures.

For a centered cyclic displacement $\mathbf m$, the normalized finite transform is
\begin{equation}
 T^{(6)}_{\mathbf m}=\frac{1}{216}
 \sum_{\kvec\in\mathcal K_6}e^{-2\pi i\kvec\cdot\mathbf m}H(\kvec),
 \qquad
 H^{(6)}(\kvec)=\sum_{\mathbf m}e^{2\pi i\kvec\cdot\mathbf m}T^{(6)}_{\mathbf m}.
 \label{eq:finite-cyclic-transform}
\end{equation}
Construction-point reconstruction, block Hermiticity, and Parseval consistency are
checked.  Because an even finite grid contains Nyquist representatives whose ordinary
integer conjugate is not represented independently, cyclic blocks are then converted
explicitly to an inversion-complete ordinary inventory.  Conversion-removed power is
recorded separately.  The retained ordinary components are
$m_i\in\{-2,-1,0,1,2\}$, giving 125 blocks before support truncation.

Complete supports are defined by the FCC metric.  For
$\mathbf m=(n_1,n_2,n_3)$,
\begin{equation}
 q(\mathbf m)=2(n_1^2+n_2^2+n_3^2+n_1n_2+n_1n_3+n_2n_3),
 \qquad
 r(\mathbf m)=\sqrt{\frac{q(\mathbf m)}{4/a^2}}.
 \label{eq:fcc-support-metric}
\end{equation}
The formula follows directly from $\Rvec=A\mathbf m$.  For the FCC primitive basis in
Eq.~\eqref{eq:silicon-direct-basis},
\begin{equation}
 \norm{A\mathbf m}^2=\frac{a^2}{2}
 (n_1^2+n_2^2+n_3^2+n_1n_2+n_1n_3+n_2n_3)
 =\frac{a^2}{4}q(\mathbf m),
 \label{eq:fcc-radius-derivation}
\end{equation}
which is equivalent to Eq.~\eqref{eq:fcc-support-metric}.  The attained labels are
ordered up to $q=48$; every support contains all blocks at or below one attained label
and is inversion complete.  These are distances between lattice-block representatives.
They are not orbital-center or individual matrix-entry distances.

For each representation and support, the producer retains maximum and RMS matrix
Frobenius errors, maximum and RMS sorted-spectrum errors, omitted block power, physical
radius, displacement inventory, and block count.  The matrix RMS is normalized by
$\sqrt{10}\times1$~eV and the sorted-spectrum RMS by the same per-state scale.  The
calculation distinguishes:
\begin{enumerate}[leftmargin=2em]
 \item \emph{direct-parent error}: the complete Wannier source relative to revealed DFT
 eigenvalues, without defining a DFT operator subtraction in the Wannier basis;
 \item \emph{exact gauge equivalence}: the complete source and attacked fields before
 finite Fourier representation;
 \item \emph{conversion}: centered cyclic blocks removed or remapped by the ordinary
 lattice convention;
 \item \emph{interpolation}: the full 125-block ordinary interpolant relative to the
 direct native represented target on the 408 points;
 \item \emph{deliberate truncation}: a support-limited ordinary model relative to its
 full ordinary interpolant; and
 \item \emph{represented total}: the final support-limited model relative to the direct
 native represented source.
\end{enumerate}
Separately maximized or RMS channels are never added as if they were signed scalar
errors.  Complete support curves are compared by observed discrete non-dominance and by
``equal-attained-error'' queries: for an error actually attained by one discrete support,
the minimum support of the other representation attaining no larger error is reported.
No interpolation between support counts and no post-hoc acceptance threshold is used.
Native Wigner--Seitz shells remain a separate fixed-gauge route.

\subsection{Native shell hierarchy and literal reduction}

The effective native representatives are grouped by Cartesian radius after rounding the
radius to nine decimal places.  This produces 18 nested, inversion-complete shells.
Table~\ref{tab:native-shells} gives the complete family.  H-shell-00 is the onsite
control and H-shell-17 is the full 383-representative source.  Literal reductions
H-shell-00 through H-shell-16 retain the source coefficients unchanged inside one shell
boundary and set every coefficient outside it to zero.  Full support remains a zero-loss
control and is excluded from proper-reduction feasibility.

\begin{table}[tbp]
\centering
\small
\caption{Complete native Wigner--Seitz shell hierarchy.  The last column is cumulative
support.}
\label{tab:native-shells}
\begin{tabular}{@{}lrrr@{}}
\toprule
Shell & Radius (\AA) & New representatives & Cumulative representatives\\
\midrule
H-shell-00&0.000000000&1&1\\
H-shell-01&3.867295746&12&13\\
H-shell-02&5.469182094&6&19\\
H-shell-03&6.698352720&24&43\\
H-shell-04&7.734591492&12&55\\
H-shell-05&8.647536177&24&79\\
H-shell-06&9.472901262&8&87\\
H-shell-07&10.231902790&48&135\\
H-shell-08&10.938364187&6&141\\
H-shell-09&11.601887238&36&177\\
H-shell-10&12.229462943&24&201\\
H-shell-11&12.826368943&24&225\\
H-shell-12&13.396705440&24&249\\
H-shell-13&13.943733109&64&313\\
H-shell-14&14.977972019&40&353\\
H-shell-15&15.945268846&12&365\\
H-shell-16&16.407546281&6&371\\
H-shell-17&16.857151341&12&383\\
\bottomrule
\end{tabular}
\end{table}

For every literal candidate, $H_S(\kvec)$ is reconstructed directly on the ordered
408-point inventory in the common Wannier basis.  The losses are
Eqs.~\eqref{eq:tau1-main}--\eqref{eq:tau2-main}; no DFT parent values are consumed.
Maximum ordered-spectrum and per-state Frobenius errors are retained as additional
coordinates.  A threshold cell $(t_1,t_2)$ is populated by the smallest proper support
with $\tau_1\le t_1$ and $\tau_2\le t_2$.  Because the cells are evaluated only after
all candidates are known, the resulting surface is descriptive rather than a
prospective selection policy.

\subsection{Refitted complete-shell family and analytic gradients}

The refitted family retains the same supports and Wannier basis but assigns one bounded
real scale to each complete retained shell.  Define the shell-resolved field
\begin{equation}
 H_s(\kvec)=\sum_{\Rvec\in s}e^{i\kvec\cdot\Rvec}T_{\rm W}(\Rvec).
 \label{eq:shell-field}
\end{equation}
For support $S_p=\{0,\ldots,p\}$,
\begin{equation}
 H_{\boldsymbol\alpha}^{(p)}(\kvec)
 =\sum_{s=0}^{p}\alpha_sH_s(\kvec),
 \qquad 0\le\alpha_s\le2.
 \label{eq:shell-refit}
\end{equation}
The same $\alpha_s$ applies to every representative in shell $s$, including both
$\Rvec$ and $-\Rvec$.  Blocks outside $S_p$ remain zero.  Sign reversal and
more-than-doubling are excluded.  No matrix entry is fitted independently, and no
atomic-orbital or Slater--Koster semantics are introduced.

Let $D_j=H_{\boldsymbol\alpha}(\kvec_j)-H_{\rm W}(\kvec_j)$ and let
$\delta\epsilon_{jn}$ be the ordered eigenvalue difference.  With $N_k=408$, rank
$N=10$, and $E_0=1$~eV,
\begin{align}
 \tau_2^2&=\frac{1}{N_kNE_0^2}\sum_j\norm{D_j}_F^2,\\
 \frac{\partial\tau_2^2}{\partial\alpha_s}
 &=\frac{2}{N_kNE_0^2}\sum_j
 \operatorname{Re}\operatorname{tr}[D_j^\dagger H_s(\kvec_j)],
 \label{eq:operator-gradient}\
 \tau_1^2&=\frac{1}{N_kN E_0^2}\sum_{j,n}\delta\epsilon_{jn}^2,\\
 \frac{\partial\tau_1^2}{\partial\alpha_s}
 &=\frac{2}{N_kN E_0^2}\sum_{j,n}\delta\epsilon_{jn}
 \langle v_{jn}|H_s(\kvec_j)|v_{jn}\rangle.
 \label{eq:spectral-gradient}
\end{align}
Equation~\eqref{eq:spectral-gradient} is the local Hellmann--Feynman derivative for the
computed ordered eigensystem.  The ordered spectral objective is nonsmooth at exact
crossings or degeneracies; the retained gradient is therefore a local binary64
derivative, not a proof of global differentiability.

For each support and scalarization weight $w$, the optimized objective is
\begin{equation}
 F_w(\boldsymbol\alpha)=w\tau_1^2+(1-w)\tau_2^2.
 \label{eq:scalarized-loss}
\end{equation}
The 17 weights are
\begin{equation}
 w\in\{0,0.001,0.003,0.01,0.03,0.1,0.2,0.35,0.5,0.65,
 0.8,0.9,0.97,0.99,0.997,0.999,1\}.
\end{equation}
Bounded L-BFGS-B uses at most 300 iterations, gradient tolerance $10^{-10}$, and
function tolerance $10^{-14}$.  Three deterministic starts are used: the literal
$\boldsymbol\alpha=\mathbf1$ candidate, an alternating 0.75/1.25 vector, and continuation
from the preceding scalarization optimum.  The analytic gradient was checked by finite
differences; the maximum absolute discrepancy was $9.7361\times10^{-10}$ and the
maximum scaled discrepancy was $4.0729\times10^{-11}$.  Only optimizer-success records enter the final witness surface.  Among the 444 retained
successful records, 50 have one coefficient active at the lower bound and 35 have one
coefficient active at the upper bound; no retained record has more than one active lower
or upper coefficient.  Bound activity occurs only on H-shell-02 (17 upper-bound
records), H-shell-06 (50 lower-bound records), and H-shell-08 (18 upper-bound records).
The bounds therefore shape parts of the observed witness set and remain part of the
model definition.  The 444 records are constructive witnesses; a missing witness is not
a global infeasibility certificate.

This is a shell-renormalized Wannier family, not a Slater--Koster model.  The separate
Slater--Koster path stopped at a rank-9 common-space gate at Gamma and one additional
construction point and produced no fitted parameter vector.

\subsection{Frozen fidelity ladder}

Three converged witnesses were selected from post-hoc threshold cells and frozen before
broader evaluation: H-shell-05 with 79 representatives from cell $(0.1,0.15)$,
H-shell-10 with 201 representatives from $(0.02,0.04)$, and H-shell-16 with 371
representatives from $(0.001,0.002)$.  Their exact shell-scale vectors, ordered from the
onsite shell outward, are
\begin{align}
 \boldsymbol\alpha_{05}={}&(0.99939463,1.02717894,1.15464602,
 0.96389478,0.87025477,0.84340786),\\
 \boldsymbol\alpha_{10}={}&(0.99931068,1.00108654,1.01987595,
 1.04502720,0.98170703,0.92902579,0.70419673,\nonumber\\
 &\hspace{17mm}1.04750727,0.57744398,0.99055868,1.15302212),\\
 \boldsymbol\alpha_{16}={}&(1.00002718,1.00000490,1.00195671,
 1.00018218,0.99843025,0.99840401,0.99896498,\nonumber\\
 &\hspace{10mm}0.99950810,0.99768267,1.00065532,0.99738969,
 0.99854754,1.00013810,1.00004865,\nonumber\\
 &\hspace{10mm}1.01699983,1.04338006,1.00871625).
 \label{eq:frozen-alpha-vectors}
\end{align}
All later calculations evaluate these coefficients without refitting.

Table~\ref{tab:data-roles} separates construction, fitting, selection, and evaluation
inventories.  In particular, the same 408 coordinates serve as the post-hoc shell-loss
inventory, but their DFT parent eigenvalues are not used to fit shell scales or select
the three frozen witnesses.

\begin{table}[tbp]
\centering
\small
\caption{Data and reciprocal inventories by role.  ``Excluded'' means unavailable to
model construction, fitting, selection, and the reported development evaluation.}
\label{tab:data-roles}
\begin{tabularx}{\textwidth}{@{}p{0.20\textwidth}XXX@{}}
\toprule
Inventory & Construction/fitting role & Selection role & Evaluation role\\
\midrule
Relaxed cell, K10 SCF, C6-B20 NSCF
& Defines the parent occurrence and Wannier inputs
& None in the post-hoc reduction
& Provenance and source-adequacy context\\
Gamma-origin $6^3$ (216 points)
& Wannier construction and finite-cyclic gauge transform
& Attack frozen before silicon support outcomes
& Construction reconstruction and Parseval checks\\
Ordered 408-point geometry
& Shell loss and scalarized refit inventory; uniform weight $1/408$
& Post-hoc threshold cells identify frozen witnesses; no DFT parent values used
& Gauge support curves and source-relative fitted-inventory losses\\
Gamma-origin $12^3$, $20^3$, $32^3$
& None
& None; models already frozen
& Dense spectra, sampled operator RMS, and mesh stability\\
$32^3$ signed-axis sectors plus local searches
& None
& None; models already frozen
& Operational conduction sector minima and curvature diagnostics\\
Parent $L$--$\Gamma$--$X$--$\Gamma$ path (241 points)
& None; fixed-density parent visualization after model freezing
& None
& Contextual DFT/source band diagram; not a fitting or global-extremum inventory\\
Frozen 402-point confirmation set
& Excluded
& Excluded
& Excluded; values remain unobserved\\
\bottomrule
\end{tabularx}
\end{table}

The ladder has three evaluation levels.  First, exact normalized-Brillouin-zone operator
RMS is computed from Eq.~\eqref{eq:parseval-rms} using the unique effective native
coefficients.  Second, spectra and sampled operator residuals are evaluated on independent
Gamma-origin tensor products
\begin{equation}
 \mathcal G_N=\{(i/N,j/N,k/N):i,j,k=0,\ldots,N-1\},
 \qquad N\in\{12,20,32\}.
 \label{eq:dense-meshes}
\end{equation}
The retained metrics are all-band RMS, maximum, and 95th-percentile absolute ordered
spectral errors; lowest-conduction-band errors; and sampled per-state operator RMS.  Mesh
comparisons are made after the models are frozen and are therefore interpolation tests,
not additional fitting data.

Third, six signed $\Delta$-family sector minima are identified operationally.  The
$32^3$ mesh is partitioned
by dominant signed Cartesian reciprocal axis.  For each of $\pm x$, $\pm y$, and
$\pm z$, the eight lowest coarse points in the corresponding sector are retained as
starts, together with the source valley coordinate.  From each start, band 4 in zero-based
ordering is minimized by L-BFGS-B over a Cartesian displacement box
$[-0.12,0.12]^3$~\AA$^{-1}$, with at most 240 iterations, function tolerance
$10^{-14}$, and gradient tolerance $10^{-10}$.  Periodic reduced coordinates are
reconstructed using $B^{-1}$, and a candidate is retained only if its final dominant
signed-axis label matches the requested valley identity.  The lowest successful matching
minimum is selected as the identified sector minimum; a bound-activity flag is retained.
This bounded multistart procedure does not prove that the global minimum of the entire
sector has been found.  The eight-start sector search
replaced a preliminary source-centered $0.08$~\AA$^{-1}$ search that reached its bound
for the H-shell-05 $y$ valleys.

At each identified sector minimum, the primary Cartesian Hessian uses step
$h=0.002$~\AA$^{-1}$.  Diagonal and mixed entries are
\begin{align}
 K_{ii}&=\frac{E(h\mathbf e_i)-2E(0)+E(-h\mathbf e_i)}{h^2},\\
 K_{ij}&=\frac{E(h\mathbf e_i+h\mathbf e_j)-E(h\mathbf e_i-h\mathbf e_j)
 -E(-h\mathbf e_i+h\mathbf e_j)+E(-h\mathbf e_i-h\mathbf e_j)}{4h^2}.
 \label{eq:finite-hessian}
\end{align}
Diagonalization gives principal curvatures and axes.  Principal masses in electron-mass
units are $(\hbar^2/m_e)/\kappa_i$ after conversion from eV~\AA$^2$.  Gap, valley
position, Hessian, and mass differences are always reduced-model differences from the
full Wannier source computed by the identical algorithm.

\subsection{Comprehensive conduction analysis}

The final identified sector-minimum coordinates from the frozen ladder are reused
without reoptimization.
For every model and valley, Cartesian gradients and Hessians are recomputed with
$h\in\{0.001,0.002,0.005,0.01\}$~\AA$^{-1}$ using
Eq.~\eqref{eq:finite-hessian}; $0.002$~\AA$^{-1}$ is primary.  With
$C=\hbar^2/(m_e\,{
m eV})\times10^{20}$ in eV~\AA$^2$, the dimensionless inverse-mass
and mass tensors are
\begin{equation}
 M^{-1}=K/C,\qquad M=C K^{-1}.
 \label{eq:mass-tensors}
\end{equation}
If $m_{v1},m_{v2},m_{v3}$ are a valley's principal masses, the reported one-valley
summary conventions are
\begin{equation}
 m_{{\rm DOS},v}=(m_{v1}m_{v2}m_{v3})^{1/3},\qquad
 m_{{\rm cond},v}=\frac{3}{m_{v1}^{-1}+m_{v2}^{-1}+m_{v3}^{-1}}.
 \label{eq:single-valley-masses}
\end{equation}
For six valleys,
\begin{align}
 m_{{\rm DOS},6}&=\left[\sum_v(m_{v1}m_{v2}m_{v3})^{1/2}\right]^{2/3},\\
 \overline{M^{-1}}&=\frac{1}{6}\sum_v M_v^{-1},\qquad
 m_{{\rm cond},6}=\frac{3}{\operatorname{tr}\overline{M^{-1}}}.
 \label{eq:six-valley-masses}
\end{align}
These are declared geometric conventions derived from represented curvatures.  They do
not include occupations, scattering, temperature, carrier density, or a transport
solution.  Gradient residuals and mass variation over all four stencils are retained to
separate reduction effects from finite-difference sensitivity.

\subsection{Gamma valence cluster and ordered-branch diagnostic}

At Gamma, zero-based bands 1--3 define the upper valence cluster.  If $X_{\rm W}$ and
$X_{\rm r}$ are orthonormal source and candidate cluster frames, the singular values
$\sigma_i$ of $X_{\rm W}^\dagger X_{\rm r}$ give principal angles
$\theta_i=\arccos\sigma_i$.  The projector discrepancy is
$\norm{X_{\rm r}X_{\rm r}^\dagger-X_{\rm W}X_{\rm W}^\dagger}_F$.  Cluster energies,
centroid, splitting, singular values, projector distance, and principal angles are all
reported.  These quantities separate preservation of the three-dimensional subspace
from preservation of the operator within that subspace.

As an explicit diagnostic of why scalar branch masses fail, the three ordered cluster
eigenvalues are sampled along [100], [110], and [111].  Five positive points and their
negative partners are used inside each half-window
$Q\in\{0.005,0.01,0.02\}$~\AA$^{-1}$.  Each ordered branch is fitted to
\begin{equation}
 E_n(q)=c_n+a_nq^2+b_nq^4,
 \label{eq:ordered-quartic}
\end{equation}
with apparent curvature $2a_n$ and apparent hole mass
$-C/(2a_n)$ only when $2a_n<0$.  Missing masses correspond to positive fitted curvature;
they are not clipped.  Because sorting can exchange branch identity at a degeneracy,
these fits are retained solely as instability diagnostics.

\subsection{Transported three-band quadratic Hamiltonian}

The primary valence object never assigns an individual band at Gamma.  For each model,
let $X_0$ be the Gamma cluster frame.  At displaced momentum $\qvec$, diagonalization
gives cluster eigenvalues $\boldsymbol\epsilon(\qvec)$ and frame $X(\qvec)$.  Write
\begin{equation}
 X_0^\dagger X(\qvec)=L(\qvec)\Sigma(\qvec)R(\qvec)^\dagger,
 \qquad Q(\qvec)=L(\qvec)R(\qvec)^\dagger.
 \label{eq:cluster-polar}
\end{equation}
The represented displaced cluster Hamiltonian in Gamma coordinates is
\begin{equation}
 \mathcal H(\qvec)=Q(\qvec)\operatorname{diag}
 [\boldsymbol\epsilon(\qvec)]Q(\qvec)^\dagger.
 \label{eq:transported-cluster}
\end{equation}
The largest principal angle encountered during sampling is retained as a transport
condition diagnostic.  The even response
\begin{equation}
 \Delta\mathcal H_{+}(\qvec)=\frac{\mathcal H(\qvec)+\mathcal H(-\qvec)}{2}
 -\mathcal H(0)
 \label{eq:even-cluster-field}
\end{equation}
is fitted by complex linear least squares to
\begin{align}
 H_{\rm eff}(\qvec)=H_{\rm eff}(0)
 &+q_x^2B_{xx}+q_y^2B_{yy}+q_z^2B_{zz}\nonumber\\
 &+2q_xq_yB_{xy}+2q_xq_zB_{xz}+2q_yq_zB_{yz}.
 \label{eq:valence-effective}
\end{align}
The sample directions are the 13 unoriented rays obtained from all nonzero vectors in
$\{-1,0,1\}^3$ after identifying $\mathbf n$ and $-\mathbf n$.  For each half-window
$Q\in\{0.005,0.01,0.02\}$~\AA$^{-1}$, four radii
$Q/4,Q/2,3Q/4,Q$ are used on every ray.  Thus each fit uses 52 even matrix observations.
Each of the six fitted $3\times3$ matrices is explicitly Hermitian symmetrized.  The RMS
fit residual is
\begin{equation}
 \left[\frac{1}{3N_q}\sum_{a=1}^{N_q}
 \norm{\Delta\mathcal H_{+}(\qvec_a)-\Delta\mathcal H_{\rm fit}(\qvec_a)}_F^2
 \right]^{1/2}.
 \label{eq:multiband-fit-residual}
\end{equation}

For a unit direction $\mathbf n$, define
\begin{equation}
 B(\mathbf n)=n_x^2B_{xx}+n_y^2B_{yy}+n_z^2B_{zz}
 +2n_xn_yB_{xy}+2n_xn_zB_{xz}+2n_yn_zB_{yz}.
 \label{eq:direction-contraction}
\end{equation}
The curvature channels are the three eigenvalues of $2B(\mathbf n)$.  Negative channels
are converted to apparent non-SOC hole masses by $m_h/m_e=-C/\kappa$.  For source-relative
coefficient comparison, the candidate Gamma frame is first polar-aligned to the source
Gamma frame; the same unitary conjugates all six candidate $B_{ij}$.  The reported joint
coefficient error is
\begin{equation}
 \eta_B=\frac{
 [\sum_{ij}\norm{B^{\rm r}_{ij}-B^{\rm W}_{ij}}_F^2]^{1/2}}
 {[\sum_{ij}\norm{B^{\rm W}_{ij}}_F^2]^{1/2}}.
 \label{eq:quadratic-coefficient-error}
\end{equation}
Window sensitivity compares each coefficient field with the primary
$Q=0.01$~\AA$^{-1}$ field by the same normalized joint norm.

Finally, direct gaps are evaluated at Gamma, six conventional X points
$\pm(2\pi/a)\mathbf e_i$, and eight conventional L points
$(\pi/a)(\pm1,\pm1,\pm1)$ after conversion to reduced primitive coordinates.  The
minimum, maximum, and spread over equivalent X or L points are retained.  A separate
Gamma-origin $32^3$ mesh supplies every ordered-band minimum, maximum, bandwidth, VBM,
CBM, and indirect gap for all four frozen models.

\resultfigure
 {../../../../../calculations/Si/bulk/analysis/c6-b20-w21-gauge-alignment-locality/figures/experiment-schematic.pdf}
 {Schematic of the rank-10 silicon gauge and finite-support experiment.}
 {The source, internal attack, alignment channels, finite Fourier route, and support
 decisions are separate stages with separate error ownership.}
 {fig:si-experiment-schematic}

\section{Silicon results}

\subsection{Complete spectrum is invariant while locality changes}

The maximum complete-spectrum difference under the frozen attack is
$3.1974\times10^{-14}$ eV.  Identity and one-global-unitary alignment leave maximum
frame defects near 1.021 and 1.020.  Pointwise polar and known-inverse controls reduce
the defect to approximately $1.29\times10^{-15}$ and $5.06\times10^{-16}$.

\resultfigure
 {../../../../../calculations/Si/bulk/analysis/c6-b20-w21-gauge-alignment-locality/figures/alignment-invariance.pdf}
 {Alignment and complete-spectrum invariance for the rank-10 silicon attack.}
 {The nonzero identity/global residual is caused by the restricted alignment family;
 the complete fields remain related by the generating momentum-dependent unitary.}
 {fig:si-alignment}

The reference and identity-attacked non-onsite finite-block powers are 20.7574 and
21.5057 eV$^2$ per retained state.  The finite-cyclic interpolation maximum sorted-band
errors are approximately 0.3470 and 0.3867 eV.  These interpolation floors are not
attack-induced complete-spectrum errors; they belong to the finite construction and
conversion route.

\resultfigure
 {../../../../../calculations/Si/bulk/analysis/c6-b20-w21-gauge-alignment-locality/figures/fourier-power-redistribution.pdf}
 {Rank-10 Fourier-power redistribution under a spectrum-preserving internal attack.}
 {Block power is plotted against physical radius, with conversion-removed power retained
 as a separate channel.}
 {fig:si-power}

\resultfigure
 {../../../../../calculations/Si/bulk/analysis/c6-b20-w21-gauge-alignment-locality/figures/complete-spectrum-equivalence.pdf}
 {Complete reference and attacked spectra on the ordered 408-point inventory.}
 {The horizontal coordinate is an ordered inventory, not an interpolated band path.}
 {fig:si-spectrum}

\resultfigure
 {../../../../../calculations/Si/bulk/analysis/c6-b20-w21-gauge-alignment-locality/figures/direct-parent-overlay.pdf}
 {Direct-parent overlay for the represented source and gauge-equivalent comparison.}
 {The displayed parent discrepancy belongs to source adequacy and remains distinct from
 the exactly spectrum-preserving internal attack.}
 {fig:si-parent-overlay}

\resultfigure
 {../../../../../calculations/Si/bulk/analysis/c6-b20-w21-gauge-alignment-locality/figures/displacement-block-power.pdf}
 {Displacement-resolved block power in the converted finite representation.}
 {Individual block norms reveal redistribution that can be hidden by a single total
 non-onsite-power scalar.}
 {fig:si-displacement-power}

The observed non-dominated support-error sets change between representations.  This is the higher-rank
material-derived counterpart of the rank-two range-inflation mechanism: complete
spectral equivalence does not fix the finite support needed under a restricted coordinate
identification.

\resultfigure
 {../../../../../calculations/Si/bulk/analysis/c6-b20-w21-gauge-alignment-locality/figures/pareto-support-error-frontiers.pdf}
 {Observed discrete rank-10 non-dominated support-error sets.}
 {Markers identify the non-dominated complete supports found on the exhaustive declared
 support inventory.  Lines guide the eye
 between discrete shell choices and do not define fractional supports.}
 {fig:si-pareto}

\resultfigure
 {../../../../../calculations/Si/bulk/analysis/c6-b20-w21-gauge-alignment-locality/figures/support-error-frontiers.pdf}
 {Complete support-error curves before non-dominated-set extraction.}
 {The curves retain dominated supports and therefore expose nonmonotonic details that a
 non-dominated-only summary removes.}
 {fig:si-support-curves}

\resultfigure
 {../../../../../calculations/Si/bulk/analysis/c6-b20-w21-gauge-alignment-locality/figures/equal-attained-error-support.pdf}
 {Equal-attained-error support comparison for the two rank-10 representations.}
 {The comparison uses attained discrete errors and avoids interpreting line segments as
 continuously adjustable support.}
 {fig:si-equal-error}

\resultfigure
 {../../../../../calculations/Si/bulk/analysis/c6-b20-w21-gauge-alignment-locality/figures/transport-diagnostics.pdf}
 {Transport and sewing diagnostics for the rank-10 alignment channels.}
 {Closure and plaquette phases are numerical sewing diagnostics and are not assigned a
 topological interpretation.}
 {fig:si-transport}

\subsection{Literal shell truncation has no global two-loss conflict}

Among the 17 proper literal reductions, H-shell-16 minimizes both primary RMS losses:
\begin{equation}
 (\tau_1,\tau_2)=
 (4.4636\times10^{-4},1.5173\times10^{-3}),
\end{equation}
while retaining 371 representatives.  The tightest literal family therefore has no
global spectral--operator conflict: the same proper support minimizes both losses.
This does not show that the losses are generally interchangeable.  The family is nested,
and both coordinates are driven by the same omitted matrix field.

The post-hoc pair $(0.1,0.15)$ first admits 135 representatives, $(0.01,0.02)$ first
admits 313, and $(0.001,0.002)$ first admits 371.  The pair $(0.0004,0.001)$ has no proper
literal witness.

\resultfigure
 {../../../../../calculations/Si/bulk/analysis/c6-b20-w21-tau1-tau2-posthoc-exploration/figures-v3/single-witness-surface.pdf}
 {Minimum-support surface for literal native-shell deletion.}
 {The full 383-representative source is excluded from reduced-family feasibility.
 Threshold coordinates are descriptive post-hoc readings, not acceptance limits.}
 {fig:literal-surface}

\subsection{Refitting creates genuine spectral--operator arcs}

The bounded refitted search retains 444 converged constructive witnesses.  Fixed
supports now trace distinct spectral--operator arcs.  Observed minimum support changes
include
\begin{center}
\begin{tabular}{@{}lrr@{}}
\toprule
$(\tau_1,\tau_2)$ & Literal representatives & Refitted representatives\\
\midrule
$(1.0,1.0)$ & 19 & 13\\
$(0.1,0.15)$ & 135 & 79\\
$(0.05,0.1)$ & 177 & 135\\
$(0.02,0.04)$ & 225 & 201\\
$(0.01,0.02)$ & 313 & 313\\
$(0.001,0.002)$ & 371 & 371\\
\bottomrule
\end{tabular}
\end{center}

Within H-shell-16, an operator-oriented witness reaches approximately
$(4.04\times10^{-4},1.394\times10^{-3})$, while a spectral-oriented witness reaches
$(2.98\times10^{-4},1.745\times10^{-3})$.  The same support therefore supports a real
tradeoff after refitting.  The optimization remains constructive; gray cells do not
certify infeasibility.

\resultfigure
 {../../../../../calculations/Si/bulk/analysis/c6-b20-w21-refitted-ws-shell-witness-posthoc/figures-v2/single-refitted-witness-surface.pdf}
 {Minimum-support surface over converged refitted shell witnesses.}
 {White circles mark active fitted witnesses and red crosses mark literal shell models.
 Gray cells mean no converged witness was found by the bounded search.}
 {fig:refit-surface}

\subsection{Fitted-point errors understate broader behavior}

The frozen fidelity ladder gives the following comparison:
\begin{center}
\small
\begin{tabular}{@{}lrrrrrr@{}}
\toprule
&\multicolumn{3}{c}{Spectral RMS (eV)}&\multicolumn{3}{c}{Operator RMS (eV)}\\
\cmidrule(lr){2-4}\cmidrule(l){5-7}
Model&408&$32^3$&Ratio&408&Parseval&Ratio\\
\midrule
H-shell-05&0.07242&0.18196&2.51&0.12814&0.26004&2.03\\
H-shell-10&0.01320&0.02785&2.11&0.03229&0.05956&1.84\\
H-shell-16&0.000404&0.000657&1.62&0.001394&0.001747&1.25\\
\bottomrule
\end{tabular}
\end{center}

Sampled operator RMS agrees with exact Parseval values to at most
$5.56\times10^{-17}$ eV.  The $12^3$ to $32^3$ spectral changes remain below about
1.4\%, so the dense values are stable over the tested sequence.  The discrepancy is a
real distinction between the fitted inventory and broader evaluation.

\resultfigure
 {../../../../../calculations/Si/bulk/analysis/c6-b20-w21-wannier-only-fidelity-ladder/figures/wannier-fidelity-ladder.pdf}
 {Frozen-model fidelity ladder from fitted points to exact and dense evaluations.}
 {Supports and shell scales are unchanged at every level.}
 {fig:si-fidelity-ladder}

\subsection{Conduction valleys and masses are stricter tests}

Relative to the full Wannier source, H-shell-05 changes the indirect gap by -65.77 meV,
displaces a valley by as much as 0.09481~\AA$^{-1}$, changes a Hessian by 29.9\%, and
changes a principal mass by 67.7\%.  H-shell-10 changes the gap by -34.09 meV and a
principal mass by 53.2\%, despite a $32^3$ all-band spectral RMS below 0.028 eV.
H-shell-16 changes the gap by -0.503 meV, the maximum valley position by
0.002758~\AA$^{-1}$, the maximum Hessian by 0.259\%, and the maximum principal mass by
2.54\%.

\resultfigure
 {../../../../../calculations/Si/bulk/analysis/c6-b20-w21-wannier-only-fidelity-ladder/figures/source-relative-feature-fidelity.pdf}
 {Source-relative conduction-edge fidelity of the frozen models.}
 {The reference is the full Wannier source, not the revealed DFT parent.}
 {fig:si-feature-fidelity}

The six-valley DOS-mass convention of the source is $1.4764m_e$, and its equal-population
conductivity-mass convention is $0.25063m_e$.  H-shell-05 changes them by 16.9\% and
20.4\%; H-shell-10 by 22.9\% and 6.90\%; H-shell-16 by 0.333\% and 0.0216\%.  Across all
models, the maximum principal-mass spread over the four finite-difference steps is below
0.172\%, so stencil variation does not explain the moderate-support errors.

\resultfigure
 {../../../../../calculations/Si/bulk/analysis/c6-b20-w21-wannier-only-comprehensive-band-edge/figures/derived-electron-mass-fidelity.pdf}
 {Derived conduction-valley mass summaries and finite-difference sensitivity.}
 {The DOS and conductivity masses are declared curvature conventions and are not
 carrier-density or transport calculations.}
 {fig:electron-mass-summary}

\subsection{The valence manifold requires a multiband operator}

The source Gamma three-band splitting is only 0.02495 meV.  H-shell-05 increases it to
102.175 meV, H-shell-10 to 54.806 meV, and H-shell-16 to 2.496 meV.  The corresponding
maximum principal angles from the source cluster are only 1.856, 0.925, and 0.0420
degrees.  A candidate can therefore preserve a recognizable three-dimensional subspace
while badly damaging the operator inside that subspace.

\resultfigure
 {../../../../../calculations/Si/bulk/analysis/c6-b20-w21-wannier-only-comprehensive-band-edge/figures/gamma-valence-cluster-fidelity.pdf}
 {Gamma valence-cluster splitting and source-subspace rotation.}
 {Internal degeneracy and subspace preservation are independent fidelity conditions.}
 {fig:gamma-valence}

\resultfigure
 {../../../../../calculations/Si/bulk/analysis/c6-b20-w21-wannier-only-comprehensive-band-edge/figures/directional-valence-cluster-dispersions.pdf}
 {Directional ordered valence-cluster dispersions near Gamma.}
 {The curves are diagnostics of branch ordering and degeneracy lifting, not a physical
 assignment of heavy-hole, light-hole, and split-off bands.}
 {fig:ordered-valence-dispersions}

Independently ordered source branches do not produce stable hole masses.  The upper
ordered branch is not hole-like for every tested window in any of the three directions,
and finite ordered-branch masses can change by about 75.5\%.  H-shell-05 and H-shell-10
appear more stable only because they artificially split the manifold.

\resultfigure
 {../../../../../calculations/Si/bulk/analysis/c6-b20-w21-wannier-only-comprehensive-band-edge/figures/directional-valence-mass-sensitivity.pdf}
 {Instability of independently ordered valence-branch mass fits.}
 {Window dependence and loss of hole-like curvature motivate replacing scalar branch
 masses with a transported multiband quadratic operator.}
 {fig:ordered-valence-mass-sensitivity}

The transported quadratic operator avoids that branch assignment.  At the primary
0.01~\AA$^{-1}$ window, the source apparent channel masses are
\begin{center}
\begin{tabular}{@{}lrrr@{}}
\toprule
Direction&Channel 1&Channel 2&Channel 3\\
\midrule
$[100]$&0.1411&0.2771&0.3015\\
$[110]$&0.09690&0.3119&2.923\\
$[111]$&0.08770&0.6971&0.7722\\
\bottomrule
\end{tabular}
\end{center}
in electron-mass units.  These are non-SOC curvature-channel eigenvalues, not physical
HH, LH, or split-off masses.

After cluster alignment, H-shell-05, H-shell-10, and H-shell-16 have quadratic-coefficient
errors of 9.54\%, 5.88\%, and 0.279\%, with maximum channel-mass changes of 38.9\%,
19.7\%, and 4.53\%.  Fit-window changes in the joint coefficient field remain below
about 0.41\%, and primary fit residuals are approximately $5$--$6\times10^{-7}$ eV RMS.

\resultfigure
 {../../../../../calculations/Si/bulk/analysis/c6-b20-w21-wannier-only-comprehensive-band-edge/figures/multiband-valence-curvature-channels.pdf}
 {Degeneracy-aware non-SOC valence curvature channels.}
 {The channels are eigenvalues of the direction-contracted quadratic matrix and are not
 assigned SOC hole labels.}
 {fig:valence-channels}

\resultfigure
 {../../../../../calculations/Si/bulk/analysis/c6-b20-w21-wannier-only-comprehensive-band-edge/figures/multiband-valence-fidelity-summary.pdf}
 {Summary of valence degeneracy, quadratic-operator, and curvature consequences.}
 {The three panels separate failure modes that a single band RMS cannot represent.}
 {fig:valence-summary}

\subsection{High-symmetry equivalence and bandwidths}

The full source's direct-gap spread among nominally equivalent X points is 0.00366 meV
and among L points is 0.0239 meV.  H-shell-05 increases these spreads to 150.1 and
68.7 meV; H-shell-10 to 36.7 and 63.9 meV; H-shell-16 to 0.427 and 2.97 meV.  Moderate
shell refitting can therefore amplify small residual symmetry differences through
cancellations among retained blocks.

On the independent $32^3$ mesh, all four valence maxima remain at Gamma.  Maximum
ordered-band bandwidth changes are approximately 354 meV, 81.1 meV, and 3.14 meV for
H-shell-05, H-shell-10, and H-shell-16.

\resultfigure
 {../../../../../calculations/Si/bulk/analysis/c6-b20-w21-wannier-only-comprehensive-band-edge/figures/high-symmetry-edge-fidelity.pdf}
 {High-symmetry direct-gap and equivalent-point diagnostics.}
 {The spread characterizes the represented models; it is not an independent validation
 of exact crystal symmetry.}
 {fig:high-symmetry}

\clearpage
\section{Discussion}

\subsection{Related mechanism under distinct evidential roles}

The rank-two theorem and rank-10 silicon attack differ greatly in dimension, origin,
geometry, and numerical complexity.  The former proves an identity in a restricted
family; the latter is a source-specific computational stress test.  They nevertheless
isolate the same logical fact.  Pointwise
unitary equivalence fixes complete spectra and leaves the represented abstract operator
unchanged under the generating map.  It does not fix the matrix field under a restricted
coordinate identification, and it does not fix the Fourier block distribution.

In the flat rank-two case, the mismatch has an exact closed form and equals induced
non-onsite power.  In the dispersive and rank-10 cases, no universal equality is claimed.
Instead, block power, support curves, transport diagnostics, interpolation, and
truncation consequences are evaluated independently.  This is a stronger scientific
position than extending the exact theorem by analogy: the theorem identifies the
mechanism, and the calculations test which parts survive when its assumptions are
relaxed.

\subsection{Representation efficiency is not physical inequivalence}

The reference and attacked complete fields are physically equivalent within the fixed
retained space under the known momentum-dependent unitary.  The reported residuals arise
because the comparison admits a smaller family: one global real rotation in the
rank-two theory, or identity/global/polar/oracle channels in the rank-10 application.
Likewise, range inflation appears only after a support family and accuracy criterion are
chosen.

The phrase ``less local'' therefore means more expensive in the declared representation
under the declared finite-support approximation.  It does not mean that the complete
gauge-related Hamiltonian describes different physics.  This distinction is essential
when locality is used as a computational resource or model-selection coordinate.

\subsection{Literal deletion and refitting answer different questions}

Literal nested shell deletion asks what happens when the source coefficients are kept
unchanged inside a radius and set exactly to zero outside it.  Its near-linear
spectral--operator requirement curve is a property of a one-coordinate nested family.
The result that H-shell-16 minimizes both losses does not prove general compatibility of
spectral and operator objectives.

Refitted shells introduce additional coordinates while preserving complete native
shells and the common Wannier basis.  Genuine same-support tradeoff arcs then appear,
and moderate post-hoc support requirements improve.  The distinction is scientifically
useful: the literal study measures deletion error, while the fitted study measures what
can be compensated by bounded renormalization of retained shells.

At tight errors, both approaches still require near-full support.  Refitting does not
turn the source into a compact 10-orbital short-range model.  It supplies explicit
witnesses and a more expressive finite family, not a global optimality theorem.

\subsection{A fitted inventory is not a Brillouin-zone guarantee}

The 408-point inventory supports exact finite statements about the loss on those points.
It does not by itself support all-$\kvec$ interpolation claims.  The fidelity ladder
freezes the models and then asks broader questions.  Parseval provides exact integrated
operator RMS for the finite Fourier series.  Dense meshes test spectra between the
fitted points.  Both reveal larger errors, particularly at moderate support.

This distinction resolves a common confusion about held-out data.  Held-out points are
not required to report a finite training-inventory result honestly.  They are required
when the claim changes from ``this model fits these points'' to ``this model interpolates
the represented Hamiltonian over a broader domain.''

\subsection{Band edges are structurally more demanding than global spectra}

Conduction masses and valence curvature depend on derivatives.  A global all-band RMS
can remain moderate while a local curvature changes by tens of percent.  H-shell-10 is
the clearest example: its dense spectral RMS is below 0.028 eV, yet at least one
principal electron mass changes by about 53\% and a non-SOC valence curvature channel
changes by about 20\%.

The valence analysis adds a qualitative warning.  Independent ordered-band Hessians can
become stable only after an incorrect reduction lifts the degeneracy.  Numerical
stability is not fidelity if the model has changed the state-space structure that made
the original scalar quantity ill defined.  A matrix-valued multiband expansion is the
appropriate retained object.

\subsection{What can and cannot be said about holes}

The current source is scalar-relativistic and non-SOC.  It supports a three-band orbital
valence operator and directional curvature channels.  It does not support standard
physical HH/LH/split-off labels, a split-off energy, or final boron-acceptor modeling.
Those claims require a spinor state space and a separately qualified SOC DFT/Wannier
branch.  SOC is not required for the internal conclusion that reduction damages the
non-SOC valence manifold.

\subsection{Why the source inadequacy matters}

The strongest positive statement about H-shell-16 is that it closely preserves the full
Wannier source across the tested internal channels.  The strongest negative statement
about H-shell-05 and H-shell-10 is that they do not.  Neither statement validates the
source.  The revealed DFT-parent comparison already shows that the source misses the
declared gap and electron-mass targets.

The one-dimensional calibration hierarchy makes the limitation operational.  Exact
spectral compression of a finite matrix does not validate that matrix against its
continuum parent, and grid refinement at fixed box width converges to the finite-box
problem rather than repairing the boundary truncation.  Likewise, enlarging the hopping
support at fixed Wannier rank converges to the fixed Wannier source.  It cannot repair a
band-window, construction-mesh, parent-model, or source/DFT discrepancy upstream.

Keeping this separation explicit prevents a common error chain:
\begin{equation*}
 \text{small reduced/source error}
 \not\Rightarrow
 \text{small source/DFT error}
 \not\Rightarrow
 \text{small DFT/experiment error}.
\end{equation*}
Each arrow requires its own calculation and authority.

\subsection{Slater--Koster is not the active conclusion}

The original Paper D plan proposed compatibility inside an orthogonal ten-orbital
$sp^3s^*$ Slater--Koster hierarchy.  The attempted common-space gate found rank 9 at
Gamma and one additional construction point.  The joint operator fit therefore stopped
before optimization.  A spectral-only Slater--Koster fit remains mathematically possible,
but it would answer a narrower eigenvalue question and is not the active reduction path.

The completed program instead studies compression directly in the common rank-10
Wannier basis.  This avoids inventing an orbital alignment that the evidence does not
support and permits finite operator, gauge, Fourier, truncation, and band-edge comparisons.
The present manuscript does not claim that Slater--Koster models are impossible; it
reports that the planned joint operator route was blocked and that a Wannier-only route
was completed instead.

\subsection{Evidence maturity}

The analytic rank-two statements are mathematical results under explicit assumptions.
The synthetic non-flat surfaces are development calculations; one active-family direct
execution independently verifies a bounded subset but not the entire five-scale preview.
The silicon analyses authenticate inputs and producer outputs but lack an independent
reconstruction.  None is a sealed public archival package.

This layered status affects publication language, not the numerical values themselves.
The correct response is to label each claim by its evidence owner and scope, rather than
either suppressing useful development results or presenting them as publication-qualified
validation.

\section{Conclusions}

The combined theory and calculations establish six main conclusions.

\begin{enumerate}[leftmargin=2em]
 \item Band equivalence does not determine a matrix representation or its hopping
 locality.  Momentum-dependent unitary similarity can preserve every complete band
 while redistributing Fourier weight.
 \item In a flat rank-two real-rotation family, the restricted-alignment threshold obeys
 an exact Bessel law, and its square equals the induced non-onsite hopping power.
 \item Representation differences can become finite-range costs at equal sampled
 accuracy.  In the controlled dispersive endpoint, the first passing cutoff rises from
 4 to 10 at $10^{-8}E_{\rm ref}$.
 \item The rank-10 silicon source exhibits the same qualitative mechanism: complete
 spectra remain invariant while alignment residuals, block power, and support-error
 frontiers change.
 \item Literal shell deletion and bounded shell refitting produce different candidate
 geometries.  Refitting creates genuine spectral--operator tradeoffs, but tight
 fidelity still requires near-full support.
 \item Global fitting losses do not control dense interpolation or band-edge fidelity.
 Conduction masses, valence degeneracy, multiband curvature, and symmetry-equivalent
 gaps expose errors that global spectral RMS can hide.
\end{enumerate}

For the studied source, H-shell-16 is the only frozen reduced witness that remains close
across all tested internal channels.  H-shell-05 and H-shell-10 are not adequate
band-edge reductions.  Even H-shell-16 does not validate silicon physics because the
full source fails separate DFT-parent criteria.

A stronger future program would independently reconstruct the silicon analyses, retain
the same frozen models for a separately labeled DFT overlay, and construct an SOC
spinor source only if physical hole or acceptor claims are required.  Any prospective
acceptance threshold must be declared before a future confirmation set is inspected.

\clearpage
\appendix

\section{Derivation details for the rank-two loss geometry}

Let a traceless Hermitian $2\times2$ matrix be written as
$D=x\sigma_x+y\sigma_y+z\sigma_z$.  Conjugation by the real rotation $R(\phi)$ leaves
the $\sigma_y$ component invariant and rotates $(x,z)$ through angle $2\phi$.  The
weighted correlation with an attacked field therefore has the form
\begin{equation}
 b_\phi=b_0+c\cos2\phi+d\sin2\phi.
\end{equation}
The continuous maximum is
\begin{equation}
 b_{\max}=b_0+\sqrt{c^2+d^2},\qquad
 \mu_{\max}=b_{\max}/\kappa.
\end{equation}
When $c^2+d^2>0$, a canonical maximizing angle is
$\phi_*=\tfrac12\operatorname{atan2}(d,c)$ reduced to the half-open fundamental
interval.  When the harmonic amplitude vanishes, the optimum value remains defined but
the angle is nonunique.

The spectral residual separates into a scalar trace shift and a traceless splitting
change.  Trace and traceless matrices are Frobenius orthogonal, which gives
Eq.~\eqref{eq:spec-quadratic}.  The operator loss expands as
\begin{align}
 L_{\rm op,\phi}^2
 &=s^2+\kappa\lambda^2-2\lambda b_\phi+\kappa\\
 &=s^2+\kappa(\lambda-\mu_\phi)^2
 +\kappa(1-\mu_\phi^2).
\end{align}
Positive $\lambda$ makes maximizing $b_\phi$ equivalent to minimizing the loss.

For the flat parent, Parseval gives
\begin{equation}
 \sum_R\norm{B_R}_F^2=\int_{-1/2}^{1/2}\norm{B(q)}_F^2dq=\norm D_F^2.
\end{equation}
The best onsite real rotation matches the norm of the zero block $B_0$.  Therefore
\begin{equation}
 \sum_{R\ne0}\norm{B_R}_F^2
 =\norm D_F^2-\norm{B_0}_F^2
 =\norm D_F^2(1-\mu_{\max}^2),
\end{equation}
which yields Eq.~\eqref{eq:residual-tail-identity} after normalization.

\section{Calculation inventory and evidence scope}

Table~\ref{tab:inventory} lists the calculation owners consumed or discussed by this
synthesis.  Inclusion in this table does not transfer evidence status between owners.

\small
\begin{longtable}{@{}p{0.30\textwidth}p{0.26\textwidth}p{0.36\textwidth}@{}}
\caption{Calculation and evidence inventory.}\label{tab:inventory}\\
\toprule
Owner & Role in synthesis & Current boundary\\
\midrule
\endfirsthead
\toprule
Owner & Role in synthesis & Current boundary\\
\midrule
\endhead
Rank-two band-equivalence/locality owner
& Exact rank-two geometry, smooth attack, dispersive range inflation, separated error channels
& Analytic results plus development calculations; passing unsealed direct-v2 covers only its declared bounded aggregate.\\
Mathieu scalar-reduction owner
& Scalar Fourier-reduction precursor
& Separately retained and independently verified; not evidence for the rank-two headline.\\
PIAB1D spectral-reduction calibration
& Separates continuum, finite-matrix, eigensolver, and invariant-subspace compression roles
& Retained development calculation with typed verification and bounded formal crosswalk; not evidence for silicon accuracy.\\
QHO1D spatial-truncation calibration
& Separates finite-domain closure, grid refinement, retained dimension, and common-coordinate pullback
& Retained unsealed development calculation; observed finite-box/grid regimes are not rigorous error bounds or silicon evidence.\\
Construction-mesh recommendation
& Freezes the $C6$ construction successor and unobserved 402-point boundary
& No calculator execution; 402 parent values remain excluded.\\
Untruncated $C6$ characterization
& Authenticates native source, reconstruction, valleys, curvatures, and locality inventory
& Structurally verified source characterization; not material adequacy.\\
Revealed-408 development audit
& DFT-parent source-adequacy audit
& Valley positions pass; gap and almost all masses fail.  Development oracle only.\\
Parent-validation preflight and revealed Davidson NSCF evidence
& Prospective criteria and direct parent comparison
& Existing authorized DFT evidence; no rerun performed for reduction fitting.\\
Accepted-S1 parent band-path evidence
& Conventional $L$--$\Gamma$--$X$--$\Gamma$ DFT diagram and source overlay
& One authorized fixed-density QE 7.5 invocation; contextual path only, with no model fitting or global-extremum claim.\\
Operator qualification
& Source reconstruction and provider eligibility
& Structural/provider status, not scientific acceptance.\\
Frozen-gauge sensitivity diagnostic
& Earlier optimizer/gauge sensitivity diagnostic
& Neither compared operator accepted; supports caution about iteration-only extension.\\
Spatial-truncation design
& Freezes native complete-shell hierarchy
& Design only; selects no reduced radius.\\
Spatial-truncation evaluation
& Evaluates every literal native support
& Source-relative results; no accepted nontrivial shell.\\
Spatial-truncation disposition
& Retains full support as the formal selected truncation
& Selected $H_T=H_W$; zero selected truncation loss.\\
Gauge/alignment/locality analysis
& Rank-10 internal gauge attack, alignment, finite Fourier, locality, support curves
& Producer checked, not independently reconstructed.\\
Literal two-loss exploration
& Literal two-loss witness surface
& Post-hoc descriptive thresholds; full support excluded from reduced feasibility.\\
Refitted shell-witness analysis
& Bounded shell-scale witnesses and fitted surface
& Constructive multistart evidence; no global infeasibility certificate.\\
Wannier-only fidelity ladder
& Parseval, dense meshes, conduction valleys, gaps, Hessians, masses
& Frozen source-relative models; no DFT validation.\\
Comprehensive band-edge analysis
& Full electron tensors, valence cluster, multiband quadratic operator, direct gaps, bandwidths
& Producer checked non-SOC development analysis; no HH/LH/SO claim.\\
Slater--Koster common-space gate
& Attempted common-space gate for joint SK fitting
& Stopped at rank 9 at Gamma and one additional point; no fit or parameter vector.\\
Hopping Pareto synthesis
& Protocol-level composition of channel-specific frontiers
& Unresolved development evaluation; selects no support.\\
Historical tutorial silicon records
& Software, parsing, and bounded same-frame decomposition controls
& Not production silicon evidence and not the source of the rank-10 results.\\
\bottomrule
\end{longtable}
\normalsize

\section{Consolidated numerical tables}

\small
\subsection{Rank-two endpoint range decisions}

\begin{center}
\begin{tabular}{@{}lrrr@{}}
\toprule
Band-error target & Reference cutoff & Comparison cutoff & Inflation\\
\midrule
$10^{-6}E_{\rm ref}$ & 3 & 8 & 5\\
$10^{-8}E_{\rm ref}$ & 4 & 10 & 6\\
$10^{-10}E_{\rm ref}$ & 6 & 12 & 6\\
$10^{-12}E_{\rm ref}$ & 7 & 14 & 7\\
\bottomrule
\end{tabular}
\end{center}

\subsection{Frozen silicon fidelity ladder}

\scriptsize
\begin{center}
\begin{tabular}{@{}lrrrr@{}}
\toprule
Model & Parseval operator RMS (eV) & $32^3$ spectral RMS (eV) & Gap change (meV) & Max mass change\\
\midrule
H-shell-05 & 0.260036 & 0.181957 & -65.77 & 67.7\%\\
H-shell-10 & 0.059558 & 0.027851 & -34.09 & 53.2\%\\
H-shell-16 & 0.001747 & 0.000657 & -0.503 & 2.54\%\\
\bottomrule
\end{tabular}
\end{center}
\small

\subsection{Valence and symmetry fidelity}

\scriptsize
\begin{center}
\begin{tabular}{@{}lrrrr@{}}
\toprule
Model & Gamma splitting (meV) & Quadratic error & Max channel-mass change & X/L gap spread (meV)\\
\midrule
Full source & 0.02495 & --- & --- & 0.00366 / 0.0239\\
H-shell-05 & 102.175 & 9.54\% & 38.9\% & 150.1 / 68.7\\
H-shell-10 & 54.806 & 5.88\% & 19.7\% & 36.7 / 63.9\\
H-shell-16 & 2.496 & 0.279\% & 4.53\% & 0.427 / 2.97\\
\bottomrule
\end{tabular}
\end{center}
\normalsize

\section{Reproducible analysis procedures}

This appendix records the analysis order in sufficient detail to prevent a later
narrative rewrite from collapsing distinct operations.  It is not a replacement for the
retained requests and producer scripts.

\subsection{Rank-10 gauge analysis}

\begin{enumerate}[leftmargin=2em]
 \item Authenticate the source lattice, rank, units, construction inventory, evaluation
 inventory, and represented matrices.
 \item Evaluate the frozen rank-10 Givens circuit and its known inverse on each requested
 reciprocal point.
 \item Construct the attacked field by pointwise unitary similarity.  Verify unitarity,
 Hermiticity, trace invariance, and complete-spectrum invariance before alignment.
 \item Evaluate identity, one-global-unitary, pointwise-polar, and known-inverse
 alignments.  Keep restricted residuals distinct from oracle residuals.
 \item On the centered $6^3$ construction tensor only, compute finite-cyclic Fourier
 blocks.  Convert them explicitly to ordinary representatives and record power removed
 by conversion.
 \item Reconstruct both finite fields and evaluate interpolation losses on the ordered
 408-point inventory.  Do not attribute these finite-route errors to the exact attack.
 \item Apply the same complete support inventory to both finite representations.
 Evaluate spectral and operator consequences at every support and extract discrete
 non-dominated points without interpolating fractional supports.
\end{enumerate}

The retained reconstruction ledger is summarized in Table~\ref{tab:gauge-ledger}.  All
listed files are relative to
\path{calculations/Si/bulk/analysis/c6-b20-w21-gauge-alignment-locality/}.  The
machine-readable files, rather than the prose description, own the exact arrays and
ordering.

\begin{table}[tbp]
\centering
\small
\caption{Machine-readable reconstruction ledger for the rank-10 gauge analysis.}
\label{tab:gauge-ledger}
\begin{tabularx}{\textwidth}{@{}p{0.23\textwidth}p{0.31\textwidth}X@{}}
\toprule
Stage & Retained artifact & Auditable content\\
\midrule
Source and coordinate contract
& \path{request.json}
& Source hashes; native and centered construction coordinates; zero offset; ordered
408-point role and weights; attack identity; tolerance policy\\
Producer implementation
& \path{run_analysis.py}
& Native HR/Wigner--Seitz reconstruction; centered permutation; transform,
conversion, support, and error-channel algorithms\\
Construction and attack result
& \path{result.json}
& \path{operator}, \path{attack}, \path{alignments}, and \path{producer_checks}; complete
spectrum and Hermiticity checks\\
Ordinary displacement inventory
& \path{result.json:displacement_power}
& Integer modes, metric-shell labels, physical radii, reference and attacked block power\\
Support and Fourier channels
& \path{result.json:alignments[*].shells}
& Exact support inventories, reconstruction, conversion/interpolation/truncation
channels, and observed non-dominated records\\
Evaluation ordering
& \path{result.json:evaluation_spectra}
& Ordered point indices, fractional coordinates, source and attacked spectra, and
separately labeled parent overlay\\
\bottomrule
\end{tabularx}
\end{table}

\subsection{Literal and refitted native-shell analyses}

The literal analysis loops over H-shell-00 through H-shell-17, reconstructs each native
Wigner--Seitz model on the 408-point inventory, and computes $\tau_1$ and $\tau_2$ with
identity alignment.  Minimum-support tables are then constructed over an explicit grid
of post-hoc threshold pairs.  H-shell-17 is retained as a control but removed from the
proper-reduction feasibility surface.

For refitted shells, each support $S$ defines a coefficient vector
$\boldsymbol\alpha\in[0,2]^{|S|}$.  The producer precomputes shell-resolved matrices and
analytic derivatives of the scalarized loss.  Seventeen scalarization weights and three
deterministic starts produce at most 51 optimization attempts per support.  Converged
solutions are deduplicated in loss/parameter space, authenticated, and retained as
constructive witnesses.  Literal points remain visible as a separate family.

The gradient audit compares analytic and finite-difference derivatives over retained
coordinates.  The maximum absolute difference is $9.7361\times10^{-10}$ and the maximum
scaled difference is $4.0729\times10^{-11}$ for the audited cases.  This verifies the
implemented derivative calculation; it does not establish global optimization.

\subsection{Frozen fidelity escalation}

The fidelity ladder is intentionally one way.  Three witness identities are copied from
the completed refitted study and frozen.  Subsequent stages may evaluate but not alter
their shell scales.

\begin{enumerate}[leftmargin=2em]
 \item Recompute the 408-point spectral and operator RMS losses from the frozen models.
 \item Compute exact integrated operator RMS by Parseval from the unique Fourier
 coefficients.
 \item Evaluate sorted spectra on independent $12^3$, $20^3$, and $32^3$ meshes.
 \item Identify six conduction sector minima through an eight-start sector search and a final
 0.12~\AA$^{-1}$ local box.
 \item Evaluate gaps, valley displacement, Hessians, and principal masses without
 refitting.
 \item Promote the same models into the comprehensive band-edge analysis.
\end{enumerate}

A preliminary valley search that encountered a 0.08~\AA$^{-1}$ search boundary was not
retained as the final result.  The final request and producer record the enlarged search
and multi-start sector logic.  This change modifies the evaluation method, not the
frozen Hamiltonians.

\subsection{Transported valence quadratic analysis}

For each model and displaced momentum $\qvec$:
\begin{enumerate}[leftmargin=2em]
 \item extract the three-dimensional upper-valence eigenspace;
 \item form its overlap with the model's Gamma cluster;
 \item compute the polar unitary and transport the displaced represented matrix into the
 Gamma-cluster coordinates;
 \item average the transported $+\qvec$ and $-\qvec$ matrices to remove odd numerical
 terms;
 \item fit the six Hermitian quadratic coefficient matrices in
 Eq.~\eqref{eq:valence-effective}; and
 \item align the candidate coefficient field to the source cluster before reporting the
 joint Frobenius error.
\end{enumerate}

This algorithm never assigns an individual band identity at the degenerate point.  The
reported channel masses are derived after direction contraction and diagonalization of
the fitted quadratic matrix.

\section{Error-channel definitions and interpretation}

\begin{longtable}{@{}p{0.23\textwidth}p{0.28\textwidth}p{0.39\textwidth}@{}}
\toprule
Channel & Primary comparison & Interpretation and exclusions\\
\midrule
Direct-parent adequacy
& DFT parent versus full Wannier source
& Tests the represented source against its parent calculation.  Not a truncation or
 gauge loss.  Current gap/mass targets fail.\\
Exact gauge equivalence
& Full attacked field versus full source spectrum
& Pointwise unitary identity.  Does not imply identity/global matrix alignment or equal
 Fourier support.\\
Restricted frame residual
& Matrix fields after identity or one-global-unitary alignment
& Measures mismatch under the declared smaller alignment family.  Not physical
 inequivalence of the complete gauge-related fields.\\
Cyclic representation
& Centered finite transform versus its exact sampled field
& Finite-discrete representation channel.  Must not be silently substituted for native
 Wigner--Seitz HR data.\\
Conversion
& Finite-cyclic representatives versus ordinary lattice representatives
& Records removed or remapped block power.  Distinct from deliberate support truncation.\\
Interpolation
& Frozen finite field on evaluation points not identical to construction points
& Measures reconstruction between construction samples.  Not a gauge error.\\
Literal truncation
& Deleted native shell model versus full source
& Coefficients inside support remain unchanged.  Measures deletion only.\\
Refitted reduction
& Bounded shell-scaled model versus full source
& Includes compensation by retained shells.  Does not have Slater--Koster orbital
 semantics.\\
Dense spectral fidelity
& Frozen model versus source on $N^3$ meshes
& Broader source-relative interpolation check.  Not DFT validation.\\
Parseval operator fidelity
& Exact coefficient difference over the normalized Brillouin zone
& Integrated same-gauge operator RMS for the finite Fourier series.\\
Conduction feature fidelity
& Source versus frozen model valleys, Hessians, and masses
& Local derivative check.  Curvature masses are not transport simulations.\\
Valence multiband fidelity
& Aligned three-band quadratic coefficient fields
& Degeneracy-aware non-SOC comparison.  No HH/LH/split-off assignment.\\
Total finite-model consequence
& Final reduced finite model versus declared source
& May contain several prior stages but is not decomposed by subtracting scalar errors
 unless an explicit identity permits it.\\
\bottomrule
\end{longtable}

\section{Artifact identification and reproducibility anchors}

The following hashes identify central retained outputs used in drafting this synthesis.
They are included as audit anchors, not as publication DOIs.

\begin{longtable}{@{}p{0.42\textwidth}p{0.50\textwidth}@{}}
\toprule
Artifact & SHA-256\\
\midrule
Literal shell \path{result.json}
& \path{2933765cb26d4aee05dca69fd73ec45f2ab26a32a145ebf9e8af80db0fdaa503}\\
Literal witness surface PDF
& \path{2184e4b5b58040aaa48adfd8926384ea3fb49ee13618d5243e5fb87e3025beec}\\
Refitted shell \path{result.json}
& \path{8a65cf3fbe3b08e8f7b71a802b8400515e1a0a514efd031e14783d6973a65664}\\
Refitted witness surface PDF
& \path{cabe28a205c6dfcdd955786b096011bc98c23a0f3fbb35bea8b11e1ec840a502}\\
Fidelity ladder \path{result-v2.json}
& \path{781c29c9512c643bbf6e23a6d524a2466caaba501dc9c7323954ba7782296fdc}\\
Fidelity ladder figure
& \path{55551ae572b1f6dc4f24971e758aed1b9a720c67ae967e328d54f71e23793b09}\\
Source-relative feature figure
& \path{0520aaf5170a640880a220ab9f28eed45e579774502fa63e4f41fec16b2a153e}\\
Comprehensive band-edge result
& \path{533dfcbed4883793bb84d1a13d1a8f62f7bb7f99b02d3ab07bb3ea99f78b0764}\\
Comprehensive multiband result
& \path{bb7be7af470c521a69897fd910fe476dce40eee1382ff8e46e6040b851718ea6}\\
Comprehensive manifest
& \path{226d654c592052a688b954009cea96b7e2b4e5c7b953d4614ca02cafa0bdf614}\\
Parent band-path \path{band-data.json}
& \path{e4a4c04eabb4cdc8f0dfef2e4b175c90ab3529b7406fbc29b2aa1aee148b5bcd}\\
Parent DFT band-diagram PDF
& \path{cd96619a6c6082e63bda4ea53f36ee737ab076151a9097989cc8642f25940bcd}\\
Narrative technical report v2 PDF
& \path{72f6850bcf71d9855d33888c3ef580ca160f0836c8e08bb0622fe6f6e966ea69}\\
Rendered Paper D PDF
& \path{bf8d4c680cf489134973539e34992df90e13d86684afc0bc9143f432e0091bd5}\\
\bottomrule
\end{longtable}

The calculation manifests, not this table, remain authoritative about complete artifact
membership.  A future rebuild of this manuscript may change its PDF hash without changing
any underlying calculation result.

\section{Scope of interpretation and remaining limitations}

\subsection{Claims supported at the present scope}

\begin{enumerate}[leftmargin=2em]
 \item Exact pointwise band equivalence does not determine hopping locality.
 \item The rank-two flat-parent threshold and residual--tail identities hold under the
 displayed assumptions.
 \item The declared dispersive development calculation exhibits positive sampled range
 inflation over most of its amplitude/scale inventory.
 \item The frozen rank-10 attack preserves the complete represented spectrum while
 changing restricted alignment and finite-block locality.
 \item Literal and refitted native-shell families have different two-loss geometry.
 \item The 408-point witness losses understate broader frozen-model errors.
 \item Moderate shell models do not preserve conduction or valence band-edge structure
 relative to the full source.
\end{enumerate}

\subsection{Claims explicitly not supported}

\begin{enumerate}[leftmargin=2em]
 \item The source is not validated as an adequate silicon Hamiltonian.
 \item No experimental silicon agreement is established.
 \item No universal hopping-range law or universal spectral/operator conversion is
 claimed.
 \item No physical HH, LH, split-off, or acceptor result is claimed.
 \item No Slater--Koster incompatibility theorem is claimed; the attempted joint fit was
 blocked by a common-space rank gate.
 \item Gray witness cells are not infeasibility certificates.
 \item Development plots and producer checks are not independent publication evidence.
 \item The frozen 402-point confirmation values were neither accessed nor inferred.
\end{enumerate}

\subsection{Before any submission or promotion}

A publication-qualified version would require a frozen claim set, reviewed numerical
controls, complete-scope retained execution where applicable, independent reconstruction
of the silicon analyses, final authorship and acknowledgments, a public archival record,
and explicit authorization for submission.  The retained provenance distinguishes the
earlier 0.5226105~eV lower-energy-frozen $4^3$/12-band candidate from the present
0.6892447~eV $C6$--$B20$--$W21$ candidate; these are different Hamiltonians compared with
the same revealed DFT parent, not conflicting estimates for one source.  If physical
hole claims are retained, a separate SOC source and state-space contract are
required.  If DFT-relative overlays are added, they must remain separate from
Wannier-relative compression losses and must not be used to refit the already frozen
models.

\section*{Data and code availability}

The rank-two calculation owner is
\path{calculations/periodic1d/rank2/band-equivalence-versus-hopping-locality/}.
The silicon analysis owners are under \path{calculations/Si/bulk/analysis/}.  Each owner
retains its own requests, source identities, producers, machine-readable results,
figures, manifests, and execution records.  The comprehensive band-edge result and its
multiband amendment are under
\path{c6-b20-w21-wannier-only-comprehensive-band-edge/}.  The narrative technical report
is under \path{calculations/Si/bulk/technical-report/v2/}.

This manuscript is a synthesis surface only.  It does not duplicate or replace the
calculation records, and no numerical value is owned solely by the prose above.

\section*{Author contributions}

E.J.M.R. developed the formulation, implemented and executed the reported development
analyses, interpreted the results, and wrote the manuscript.  Authorship and contribution
statements remain subject to final review before submission.

\section*{Conflict of interest}

The author declares no conflict of interest in this working manuscript.

\printbibliography[title={References}]

\end{document}
